\documentclass[11pt]{article}
\usepackage[T1]{fontenc}
\usepackage[utf8]{inputenc}
\usepackage{lmodern,microtype}
\usepackage[a4paper,margin=29mm]{geometry}
\usepackage{amsmath,amssymb,amsthm,mathtools,booktabs,longtable,array}
\usepackage{aliascnt}
\usepackage{graphicx,float,needspace}
\usepackage{titling}
\setlength{\droptitle}{-11mm}
\usepackage[dvipsnames]{xcolor}
\definecolor{ManuscriptLinkBlue}{HTML}{0000FF}
\usepackage[colorlinks=true,linkcolor=ManuscriptLinkBlue,citecolor=ManuscriptLinkBlue,urlcolor=ManuscriptLinkBlue]{hyperref}
\usepackage[nameinlink,noabbrev]{cleveref}
\hypersetup{
  pdftitle={Largest Three-dimensional Cubes in Higher-dimensional Unit Cubes},
  pdfauthor={K.-C. Chang},
  pdfsubject={Preliminary working draft, 25 November 2026}
}
\numberwithin{equation}{section}
\newtheorem{theorem}{Theorem}[section]
\newaliascnt{proposition}{theorem}
\newtheorem{proposition}[proposition]{Proposition}
\aliascntresetthe{proposition}
\newaliascnt{lemma}{theorem}
\newtheorem{lemma}[lemma]{Lemma}
\aliascntresetthe{lemma}
\newaliascnt{corollary}{theorem}
\newtheorem{corollary}[corollary]{Corollary}
\aliascntresetthe{corollary}
\theoremstyle{definition}
\newaliascnt{definition}{theorem}
\newtheorem{definition}[definition]{Definition}
\aliascntresetthe{definition}
\theoremstyle{definition}
\newaliascnt{remark}{theorem}
\newtheorem{remark}[remark]{Remark}
\aliascntresetthe{remark}
\crefname{theorem}{Theorem}{Theorems}
\crefname{proposition}{Proposition}{Propositions}
\crefname{lemma}{Lemma}{Lemmas}
\crefname{corollary}{Corollary}{Corollaries}
\crefname{definition}{Definition}{Definitions}
\crefname{remark}{Remark}{Remarks}
\newcommand{\R}{\mathbb R}
\newcommand{\N}{\mathbb N}
\newcommand{\F}{F_3}
\newcommand{\abs}[1]{\lvert#1\rvert}
\newcommand{\norm}[1]{\lVert#1\rVert}
\DeclareMathOperator{\tr}{tr}
\DeclareMathOperator{\diag}{diag}
\DeclareMathOperator{\rank}{rank}
\DeclareMathOperator{\supp}{supp}
\newcommand{\ds}{\delta_*}
\newcommand{\Dlt}{\Delta}
\title{Largest three-dimensional cubes\\in higher-dimensional unit cubes}
\author{\normalsize\begin{minipage}{0.87\textwidth}\centering
  {\large K.-C. Chang}\\[5pt]
  \emph{Preliminary working draft. This manuscript has not yet undergone final editorial revision.}\\[4pt]
  Questions, comments, and corrections are welcome at\\[2pt]
  \href{mailto:mathdog81321@gmail.com}{\texttt{mathdog81321@gmail.com}}.
\end{minipage}}
\date{\normalsize 25 November 2026}
\begin{document}
\maketitle
\begin{abstract}
We determine the maximal side length $f(3,n)$ of a three-dimensional
cube contained in the unit $n$-cube for every $n\ge3$. The values
for $3\mid n$ and in dimension $4$ were previously known, and the value
in dimension $5$ had been reported without a published proof.
We obtain exact values in the remaining cases, including the exceptional
dimension $n=7$, and prove the recurrence
$f(3,n+3)^2=f(3,n)^2+1$ for every $n\ge8$.

The containment problem is equivalent to optimizing real matrices with
$A^{\mathsf T}A=sI_3$ and $\|a_i\|_1\le1$ for every row $a_i$.
The sharp upper bounds use multiplier equations, second-order variations,
and estimates for rows supported on two coordinates. We also establish,
for every fixed inner dimension $k\ge2$, an equivalence between eventual
stabilization of the squared optimum and the existence of maximizing
matrices with uniformly bounded numbers of rows other than signed
coordinate vectors. A sign-matrix correspondence gives exact values for
regular tetrahedra whose centroid is the origin. A Lean~4 formalization
of the main results accompanies the manuscript.
\end{abstract}
\medskip
{\setcounter{tocdepth}{1}\small\tableofcontents}
\newpage

\section{The problem and the formula}
\label{sec:main}
Let $C_d=[-1/2,1/2]^d$.  For $n\ge k$, let $f(k,n)$ be the largest
side length of a $k$-cube contained isometrically in $C_n$.  The embedded cube
may be rotated and translated.  We ask for containment, not for a
cube passing through a hole in another one.  The latter is related to
Rupert's problem (see \cite{HSW2018,SteiningerYurkevich2023}).

We write $F_k(n)=f(k,n)^2$ and set
\[
 \ds=15-4\sqrt{14},\qquad \Dlt=10-4\sqrt6.
\]
The two exceptional algebraic constants are
$\delta_1\approx0.014925$, the smaller positive root of
\[
 p_1(d)=4d^4-12d^3-67d^2-66d+1,
\]
and $\delta_2\approx0.030276$, the smallest positive root of the
irreducible polynomial
\[
 p_2(d)=d^8-8d^7-12d^6+80d^5+70d^4-200d^3-28d^2-32d+1.
\]
Their defining equations, which come from the exceptional
constructions, are given in \cref{sec:algebraic-parameters}, where
the polynomials $p_1,p_2$ are also derived.

\begin{theorem}[Exact formula]\label{thm:sharp-values}
For every integer $n\ge3$,
\begin{equation}\label{eq:main-formula}
 \F(n)=
 \begin{cases}
 m,& n=3m,\ m\ge1,\\
 1+\delta_1,&n=4,\\
 2+\delta_2,&n=7,\\
 m+\ds,&n=3m+1,\ m\ge3,\\
 m+\Dlt,&n=3m+2,\ m\ge1.
 \end{cases}
\end{equation}
In particular, $f(3,n)=\sqrt{\F(n)}$.
\end{theorem}

The case $3\mid n$ follows from a trace bound. The other two
residue classes require separate sharp upper bounds.  We give the constructions
in \cref{sec:matrix} and prove optimality in
\cref{sec:residue-one,sec:residue-two}.

\begin{corollary}[Eventual recurrence]\label{cor:eventual-recurrence}
For every integer $n\ge8$,
\begin{equation}\label{eq:recurrence}
 f(3,n+3)^2=f(3,n)^2+1.
\end{equation}
More precisely, for $n\ge3$ the equality holds if and only if
$n\notin\{4,7\}$.
\end{corollary}
\begin{proof}
Except at $n=4,7$, the formulas in \cref{thm:sharp-values}
use the same expression for $\F(n)$ and $\F(n+3)$, with the integer
parameter $m$ increased by one. Their difference is therefore $1$.
It remains to check the transitions from dimension $4$ to $7$
and from $7$ to $10$. For these two transitions, we need to know
which of the exceptional corrections $\delta_1,\delta_2$ is larger,
and how both compare with the stable correction $\ds$.
The root-isolation estimates proved later in \cref{lem:root-isolation}
give the strict ordering
\[
 \delta_1<(17/100)^2<(347/2000)^2<\delta_2
 <49/1600<\ds.
\]
The final inequality follows from
$\sqrt{14}<23951/6400$, which holds because both sides are positive
and $23951^2-14\cdot6400^2=210401>0$.
Using the formula for the three relevant dimensions gives
\[
 \F(7)-\F(4)=1+\delta_2-\delta_1>1,\qquad
 \F(10)-\F(7)=1+\ds-\delta_2>1.
\]
These are the only two failures of the recurrence.
\end{proof}

\subsection{Earlier work}
Gardner~\cite{Gardner1966} asked for $f(3,4)$.  The general problem
$f(k,n)$ appears in \cite[Problem~B4]{CFG1991}.  Shultz
\cite{Shultz2013} obtained the divisible-dimensional values, the cases
$k=2$, and $f(3,4)$.  Finch~\cite[Sec.~8.14]{Finch2003} records the
five-dimensional value attributed to Huber:
\[
 f(3,5)=\sqrt{11-4\sqrt6}=2\sqrt2-\sqrt3.
\]
An explicit construction appears in \cite{Firsching2014}, where it is
presented as a lower bound obtained with numerical optimization and
explicitly not as a proof of optimality.
Thus $f(3,4)$ was proved before this work, and the value of $f(3,5)$
had been reported. We are not aware of a published proof of the
matching upper bound for $f(3,5)$; \cref{sec:residue-two} supplies one.
Huber and Ligocki's related study of scaling in generalized Rupert
problems is listed as unpublished in \cite{HSW2018}.
No result from that manuscript is used here.
The distinction between this cube problem and related results on
simplices is discussed in \cref{sec:consequences}.

\subsection{Proof outline}
Let $A\in\mathbb R^{n\times3}$ be an optimal matrix at level
$s=F_3(n)$.  The main constraints are
\[
 A^{\mathsf T}A=sI_3,\qquad \|a_i\|_1\le1.
\]
The trace bound gives $s\le n/3$, with equality when $3\mid n$.
For the remaining dimensions, a key step is to prove that every row
of a maximizing matrix satisfies $\|a_i\|_1=1$. In particular, the
case $n\equiv2\pmod3$ requires a separate argument to exclude rows
with zero multiplier weight. We then use multipliers and second-order
variations to bound the possible supports. A row with three nonzero
coordinates is called \emph{full}, and a row with exactly two is called
a \emph{pair row}. If $F$ and $Q$ count these rows respectively and a
full row occurs, then
\[
 Q+2F\le5.
\]
The structural classification of full-row cases appears in
\cref{sec:variation} and is shared by the two residue classes.
When there are no full rows, each remaining non-pure row uses exactly
two coordinates. The multiplier equations orient the occurring
coordinate pairs into a directed acyclic graph. Quantitative bounds
on the squared entries of these pair rows then give the stable upper
bounds. Dimensions $4$ and $7$ require separate algebraic cases.
Explicit constructions in \cref{sec:matrix,prop:exceptional-frames}
provide the matching lower bounds.

% BEGIN sections/model
\section{Matrix formulation and bounds}
\label{sec:matrix}
\begin{definition}
An $n\times k$ real matrix $A$ is feasible at level $s$ if
\[
 A^{\mathsf T}A=sI_k,\qquad \norm{a_i}_1\le1\quad(1\le i\le n),
\]
where $a_i$ denotes its $i$th row.
\end{definition}

We call a feasible matrix at positive level a \emph{feasible frame}.
A \emph{maximizing frame} is feasible at the largest possible level
$F_k(n)$. The inequality $\|a_i\|_1\le1$ is the \emph{row budget}
for row $i$. Every row is subject to its own budget.
Its \emph{support} is the set of columns in which it is nonzero,
$\supp(a_i)=\{j:a_{ij}\ne0\}$.

The condition $A^{\mathsf T}A=sI_k$ will be called \emph{isotropy at
level $s$}. It means that the $k$ columns of $A$ are mutually
orthogonal and each has squared Euclidean norm $s$. In particular,
for every column $j$,
\[
 \sum_{i=1}^{n}a_{ij}^2=s.
\]
We refer to any sum of squares of matrix entries as its
\emph{squared mass}, or \emph{energy}. More precisely, for a set of
entry positions $\mathcal I\subseteq\{1,\ldots,n\}\times
\{1,\ldots,k\}$, its squared mass is
$\sum_{(i,j)\in\mathcal I}a_{ij}^2$. Thus the energy in a whole
column is exactly $s$, whereas the mass in a subset of that column
need not equal $s$.

In the three-column case, a \emph{pure row} has support of size one,
a \emph{pair row} has support of size two, and a \emph{full row}
has support of size three. A \emph{balanced pair row} has two
nonzero entries with equal absolute values. A row is
\emph{saturated} when $\|a_i\|_1=1$ and \emph{slack} when
$\|a_i\|_1<1$. We call a pair row or a full row a \emph{mixed row}.
The \emph{mixed contribution} to squared mass in column $j$ is
\[
 E_j^{\mathrm{mix}}(A)=
 \sum_{\substack{1\le i\le n\\|\supp(a_i)|\in\{2,3\}}}a_{ij}^2.
\]
For later reference, a saturated pure row is necessarily one of the
\emph{signed unit coordinate vectors} $\pm e_j$, where $e_j$ denotes
the $j$th standard basis vector.

\begin{proposition}\label{prop:matrix-model}
For integers $k\ge1$ and $n\ge k$,
\[
 F_k(n)=\max\{s:A\text{ is feasible at level }s\}.
\]
The maximum exists, and $F_k(n)\le n/k$. Equality holds precisely when
$k$ divides $n$.
\end{proposition}
\begin{proof}
An isometric copy of a cube of side $\ell$ has the form
$c+\ell U C_k$, where $U^{\mathsf T}U=I_k$.  Its projection onto the
$i$th coordinate is an interval of length $\ell\norm{u_i}_1$.
Indeed, the maximum of $u_i\cdot x$ over $x\in C_k$ equals
$\norm{u_i}_1/2$.  Containment in $C_n$ therefore requires
$\ell\norm{u_i}_1\le1$.  Conversely, these inequalities place the
centered copy $\ell U C_k$ inside $C_n$.  Hence allowing a translation
does not improve the optimum.  Taking $A=\ell U$ gives the matrix
formulation.
The feasible matrices form a closed subset of the compact product of
$n$ unit $\ell^1$ balls, and $s=k^{-1}\tr(A^{\mathsf T}A)$ is continuous.
Furthermore,
\[
 ks=\sum_i\norm{a_i}_2^2\le\sum_i\norm{a_i}_1^2\le n.
\]
If equality holds, every row has $\ell^1$-norm $1$ and only one nonzero
entry. Hence every row is some $\pm e_j$. Because the Gram matrix
has equal diagonal entries, the number of these unit rows must be
the same in each column. That number is an integer, so $k\mid n$.
Conversely, taking $n/k$ vertically concatenated copies of $I_k$
gives equality.
\end{proof}

For matrices $A\in\R^{n_1\times k}$ and $B\in\R^{n_2\times k}$,
\emph{stacking} means placing one matrix below the other to form
$\begin{pmatrix}A\\B\end{pmatrix}\in\R^{(n_1+n_2)\times k}$.
The Gram matrix of the stacked matrix is the sum
$A^{\mathsf T}A+B^{\mathsf T}B$. Thus stacking feasible matrices
at levels $s_1,s_2$ produces a feasible matrix at level $s_1+s_2$.

\begin{proposition}[Stacking]\label{prop:stacking}
For $n\ge k$,
\[
 F_k(n+k)\ge F_k(n)+1.
\]
More generally, $F_k(n_1+n_2)\ge F_k(n_1)+F_k(n_2)$ for $n_1,n_2\ge k$.
\end{proposition}
\begin{proof}
The Gram identity above shows that stacking preserves the row budgets
and adds the two levels. For the first inequality, take $B=I_k$,
which is feasible at level $1$.
\end{proof}

\begin{remark}
Stacking gives only a lower bound.  The equality in
\cref{cor:eventual-recurrence} needs the sharp upper bound.
\end{remark}

We shall use three nonnegative quantities associated with a feasible
matrix. The \emph{row loss} $L_i$ measures how much the squared
Euclidean norm of row $i$ falls short of $1$. The \emph{pair-product
mass} $P_i$ sums the absolute products of distinct entries in that
row. Their total loss is $L$:
\[
 L_i=1-\norm{a_i}_2^2,\qquad
 P_i=\sum_{j<l}\abs{a_{ij}a_{il}},\qquad L=\sum_iL_i=n-ks.
\]
Here $j,l$ range over the $k$ columns and $i$ over the $n$ rows.
Since $\norm{a_i}_1^2=\norm{a_i}_2^2+2P_i\le1$,
we have $L_i\ge2P_i$.
For each pair of columns, orthogonality means that its signed products
sum to zero. The absolute product in one row therefore cannot exceed
the sum of the corresponding products in the other rows.
Summing this observation over all column pairs yields
\begin{equation}\label{eq:cancellation}
 \sum_{h\ne i}P_h\ge P_i.
\end{equation}

\begin{lemma}[Dominant coordinates]\label{lem:dominance}
Let $k\ge2$, $n=km+1$, and let $A$ be feasible at level $m+\delta$,
where $\delta>0$. Every row has a unique coordinate whose absolute value
exceeds the sum of the absolute values of its other coordinates.
If $a$ is that absolute value and $b$ is the sum of the others, then
\[
 a-b\ge\sqrt{k\delta},\qquad
 b\le\frac{1-\sqrt{k\delta}}2.
\]
Call a row \emph{dominant in column $j$} if that unique coordinate
is $j$, and write $d(i)$ for the dominant column of row $i$.
The squared mass outside the dominant positions is, by definition,
\[
 E=\sum_{i=1}^n\ \sum_{\substack{1\le j\le k\\j\ne d(i)}}a_{ij}^2.
\]
It satisfies
\begin{equation}\label{eq:E-bound}
 (k-1)\delta\le E\le\frac{(1-\sqrt{k\delta})^2}{2}.
\end{equation}
There are $m+1$ dominant rows in one column and $m$ in each other column.
Consequently
\begin{equation}\label{eq:delta-coarse}
 \delta\le\frac1{(\sqrt{k}+\sqrt{2(k-1)})^2}.
\end{equation}
\end{lemma}
\begin{proof}
By \eqref{eq:cancellation},
$L\ge1-\norm{a_i}_2^2+2P_i$. Since $L=1-k\delta$,
\[
 \norm{a_i}_2^2-2P_i\ge k\delta.
\]
Let $a$ be a largest coordinate magnitude, and $b$ the sum of the rest.
If $a\le b$, the fact that every coordinate has magnitude at most $a$
gives $\norm{a_i}_2^2\le a(a+b)$.  Since
$\norm{a_i}_1=a+b$ and $(a+b)^2=\norm{a_i}_2^2+2P_i$,
\[
 \norm{a_i}_2^2-2P_i
 =2\norm{a_i}_2^2-(a+b)^2
 \le(a+b)(a-b)\le0,
\]
contradicting the strict lower bound $k\delta>0$.
Thus $a>b$, and this dominant coordinate is unique.  Moreover,
\[
 \norm{a_i}_2^2-2P_i\le(a-b)^2.
\]
Combining this with the same lower bound yields
$a-b\ge\sqrt{k\delta}$.  The inequality $a+b\le1$ then gives
$b\le(1-\sqrt{k\delta})/2$.

Put $d=\sqrt{k\delta}$. The conditions $a-b\ge d$ and $a+b\le1$
give $b/a\le(1-d)/(1+d)$. Thus the nondominant squared mass of this row
is at most $b^2\le[(1-d)/(1+d)]P_i$. Summation yields
\[
 E\le\frac{1-d}{1+d}\frac{L}{2}=\frac{(1-d)^2}{2}<\frac12.
\]
A column with at most $m-1$ dominant rows would have energy at most
$m-1+E<m$, which is impossible. Counting the $km+1$ rows gives the
asserted multiplicities. The $k-1$ columns with $m$ dominant rows each need at least
$\delta$ of nondominant energy, proving the lower bound in
\eqref{eq:E-bound}. Taking square roots of that inequality proves
\eqref{eq:delta-coarse}.
\end{proof}

\begin{corollary}\label{cor:two-dimensional}
For every $m\ge1$, $F_2(2m+1)\le m+1/8$.
For a three-column feasible frame with $n=3m+1$ and level $m+\delta>m$,
\[
 \delta\le7-4\sqrt3<\frac{72}{1000}.
\]
\end{corollary}
\begin{proof}
Apply \cref{lem:dominance} with $k=2$ and $k=3$.
Levels at most $m$ require no further argument.
\end{proof}

\subsection{Two explicit cores}
In this subsection, a \emph{core} is a fixed finite block of rows
used to construct a family of feasible matrices. It may be feasible
by itself, as in the five-row example below, or require additional
unit coordinate rows to make the Gram matrix a scalar multiple of
the identity. The term describes a construction, not a uniquely
determined part of an optimal matrix. The \emph{residual core} in
\cref{prop:finite-core} has a different, precise meaning: it is the
collection of rows that are not signed unit coordinate vectors.
The following five-row matrix attains the previously reported value $F_3(5)$
\cite{Finch2003}.  Another construction appears in \cite{Firsching2014}.
We use this matrix because it can be stacked with copies of $I_3$:

\begin{equation}\label{eq:five-core}
 A_5=\begin{pmatrix}
 1&0&0\\ b&a&0\\ b&-a&0\\ b&0&a\\ b&0&-a
 \end{pmatrix},\qquad
 b=\frac{\sqrt6-2}{2},\quad a=1-b.
\end{equation}
As $a,b>0$ and $a+b=1$, every row has $\ell^1$-norm $1$,
including those containing a negative entry, since
$\|(b,-a,0)\|_1=b+a$.  Moreover,
\[
 A_5^{\mathsf T}A_5
 =\diag(1+4b^2,2a^2,2a^2)=(1+\Dlt)I_3.
\]
Stacking $m-1$ identity blocks gives $\F(3m+2)\ge m+\Dlt$.

For dimensions $n=3m+1$ with $m\ge3$, the formula will have a
constant correction term. We refer to these dimensions as the
\emph{stable first-residue range}. Set
\[
 \beta=\frac{4-\sqrt{14}}2,\qquad \alpha=1-\beta.
\]
Use the four rows
\begin{equation}\label{eq:star-mixed}
 (\alpha,\beta,0),\quad(\alpha,-\beta,0),\quad
 (\alpha,0,\beta),\quad(\alpha,0,-\beta),
\end{equation}
together with $m-3$ copies of $e_1$, $m$ copies of $e_2$, and $m$ copies
of $e_3$. This gives $3m+1$ rows. Each has $\ell^1$-norm $1$, and the
resulting Gram matrix is
\[
 \diag(m-3+4\alpha^2,m+2\beta^2,m+2\beta^2)
 =(m+\ds)I_3.
\]
This is the \emph{four-pair star construction}, using only column
pairs $\{1,2\}$ and $\{1,3\}$.
Here $m\ge3$, since the number of copies of $e_1$ is $m-3$.
The case $n=7$ needs a different construction.

% END sections/model
% BEGIN sections/variational
\section{Multipliers and support}
\label{sec:variation}
\subsection{Multiplier systems and rank restrictions}
Fix an optimal feasible matrix $A$ at level $s=\F(n)>0$.
Consider the set of all $n\times3$ matrices with the same Gram
matrix as $A$:
\[
 \mathcal M_s=\{B\in\R^{n\times3}:B^{\mathsf T}B=sI_3\}.
\]
This is a scaled \emph{Stiefel manifold}. We study the function
$h(B)=\max_i\norm{b_i}_1$, the largest row $\ell^1$-norm.
It has minimum $1$ at $A$.
Indeed, $h(A)\le1$, and if any $B\in\mathcal M_s$ had $h(B)<1$,
then $B/h(B)$ would be feasible at level $s/h(B)^2>s$.
Differentiating the Gram constraint along a curve through $A$ gives
\[
 T_A\mathcal M_s=\{H:A^{\mathsf T}H+H^{\mathsf T}A=0\}.
\]
This tangent space has dimension $3n-6$. We use the Frobenius
inner product $\langle X,Y\rangle=\tr(X^{\mathsf T}Y)$.
The corresponding normal space is $\{AS:S=S^{\mathsf T}\}$, since
$A^{\mathsf T}H$ is skew-symmetric for every tangent $H$ and the
symmetric matrices form a six-dimensional space.
Finally, write $h$ as the maximum of the finitely many linear forms
$\ell_{i,\varepsilon}(B)=\varepsilon^{\mathsf T}b_i$,
$\varepsilon\in\{-1,1\}^3$. Each inequality
$\ell_{i,\varepsilon}(B)\le1$ is one of the linear \emph{facet
constraints} expressing the row budget. It is \emph{active at $A$}
precisely when $\ell_{i,\varepsilon}(A)=1$. We use \emph{active facet}
to mean an active constraint, not a separate condition on a row.

For $v\in\R^3$, write $\partial\|v\|_1$ for the \emph{subdifferential}
of the $\ell^1$-norm. In concrete terms, it consists of those
$g\in[-1,1]^3$ with $g_j=\operatorname{sgn}(v_j)$ whenever $v_j\ne0$.
At zero coordinates, any value in $[-1,1]$ is allowed.

\begin{lemma}[Normalized multipliers]\label{lem:multipliers}
There is a nonempty compact set of normalized multiplier systems
$(\lambda,g,S)$ satisfying
\begin{gather*}
 G=AS,\qquad S=S^{\mathsf T},\qquad
 G_{i,*}=\lambda_i g_i^{\mathsf T},\\
 \lambda_i\ge0,\qquad \sum_i\lambda_i=1,\qquad
 g_i\in\partial\norm{a_i}_1,
 \qquad \lambda_i>0\Longrightarrow\norm{a_i}_1=1.
\end{gather*}
The scalars $\lambda_i$ are the \emph{row multiplier weights},
and $S$ is the \emph{multiplier matrix}. Positive weight implies
a saturated row, but the converse is not assumed.
Normalization means $\sum_i\lambda_i=1$. The weights are
associated with a chosen multiplier system, not fixed attributes
of the rows of $A$. For each such system,
\begin{equation}\label{eq:multiplier-bounds}
 S_{jj}\ge0,\qquad
 \abs{S_{jk}}\le\min(S_{jj},S_{kk}),\qquad \tr S=1/s.
\end{equation}
\end{lemma}
\begin{proof}
No tangent vector can make all active forms strictly negative.
Otherwise a smooth curve in $\mathcal M_s$ with that initial velocity
would decrease every active form and, by continuity and finiteness, all
other forms would remain below $1$.  This would contradict the
minimality of $A$.

We use the following finite-dimensional alternative, which is
Gordan's theorem of the alternative (see, e.g.,
\cite[Ch.~7]{Schrijver1986}).  For linear
functionals $L_1,\ldots,L_N$ on a vector space $T$, either some
$v\in T$ satisfies $L_a(v)<0$ for every $a$, or there are
$\mu_a\ge0$ with $\sum_a\mu_a=1$ such that
$\sum_a\mu_aL_a=0$ on $T$.  These are mutually exclusive and
exhaustive, as follows by separating the image of
$v\mapsto(L_1(v),\ldots,L_N(v))$ from the open negative orthant.
Apply this alternative to the restrictions of the active forms to
$T_A\mathcal M_s$.  The resulting convex combination of their
gradients annihilates the tangent space, so it equals $AS$ for a
symmetric matrix $S$.

Group the facet weights belonging to row $i$ into $\lambda_i$.
When $\lambda_i>0$, their normalized convex combination gives
$g_i\in\partial\norm{a_i}_1$.  For a zero weight, choose any
subgradient of that row.  The conditions on the $\lambda_i$ follow
because only active facets appear.  The set of systems is compact:
facet weights lie in a closed subset of a finite simplex, subgradients
lie in closed subsets of $[-1,1]^3$, and
$S=s^{-1}A^{\mathsf T}G$ is continuous in these data.

By the explicit description of the subdifferential above,
$|g_{ij}|\le1$ and $a_{ij}g_{ij}=|a_{ij}|$. It follows that
\[
 S_{jj}=s^{-1}\sum_i\lambda_i|a_{ij}|,\qquad
 |S_{jk}|\le s^{-1}\sum_i\lambda_i|a_{ij}|=S_{jj}.
\]
Symmetry gives the other bound. Taking the trace proves the last identity.
\end{proof}

\begin{lemma}[Exclusion of rank one]\label{lem:rank-one}
Suppose $n=3m+r$, $r\in\{1,2\}$, and $s=m+\delta$.
For $r=1$, assume $0<\delta<1/3$. For $r=2$, assume
$1/5<\delta\le4/9$. Then no normalized multiplier matrix has rank one.
Every singular multiplier matrix therefore has rank two and is positive
semidefinite.
\end{lemma}
\begin{proof}
If $S$ has rank one, its positive trace gives $S=\sigma xx^{\mathsf T}$,
where $\sigma>0$ and $\norm{x}_2=1$. Let $M$ be the set of coordinates
of maximum absolute value, $q=|M|$, and let $r_+$ count the positive
row multipliers. A positive-multiplier row must be supported on $M$,
have $\ell^1$-norm $1$, and have its coordinate signs aligned with $x$
up to one common sign. Consequently
\[
 s=\norm{Ax}_2^2=r_+\norm{x}_\infty^2\le r_+/q.
\]
Apply isotropy to the unit vector equal to $\operatorname{sgn}(x_j)/\sqrt q$
on $M$ and zero elsewhere. The positive rows alone contribute $r_+/q$,
so $s\ge r_+/q$. Equality follows, and $s=r_+/q$.
In the first residue class this is incompatible with $0<\delta<1/3$.
In the second, the only possibility is $q=3$, $\delta=1/3$, and
$r_+=3m+1$. After changing coordinate signs, all products within the
positive rows are nonnegative. Only one row remains to cancel them.
That row cannot have three nonzero entries, since some pair of its
coordinates has positive product. If the remaining row is supported on two coordinates, every positive
row involving the third coordinate must be pure, since its product
with any other coordinate is nonnegative and cannot be canceled.
That column then has integer squared energy, contrary to $s=m+1/3$.
If the remaining row has support of size at most one, the same argument
applies to at least two columns.

All principal minors of order two are nonnegative by
\eqref{eq:multiplier-bounds}. If $S$ has rank two, let $\mu_1,\mu_2$
be its nonzero eigenvalues. Then
\[
 \mu_1\mu_2=\sum_{j<k}(S_{jj}S_{kk}-S_{jk}^2)\ge0,
 \qquad \mu_1+\mu_2=\tr S=\frac1s>0.
\]
Since neither eigenvalue vanishes, both are positive.
\end{proof}

\begin{lemma}[Active rows in the first residue class]\label{lem:active-r1}
If $n=3m+1$ and $\F(n)=m+\delta>m$, every row multiplier in every
normalized system is positive. In particular, every row has
$\ell^1$-norm $1$.
\end{lemma}
\begin{proof}
The bound in \cref{cor:two-dimensional} excludes rank one.
Suppose $\lambda_i=0$ for some row. By \cref{lem:dominance}, this
row is nonzero, so $S$ is singular. The rank-one exclusion gives
$\ker S=\R x$. Each zero-weight row is parallel to $x$, which has
a strictly dominant coordinate $j$:
\[
 |x_j|>\sum_{k\ne j}|x_k|.
\]
For $\lambda_i>0$, the multiplier equation gives
$g_i^{\mathsf T}x=0$. If $a_{ij}\ne0$, then $|g_{ij}|=1$, contradicting
strict dominance. Hence
\[
 \lambda_i>0\Longrightarrow a_{ij}=0,\qquad
 \lambda_i=0\Longrightarrow a_i^{\mathsf T}=t_i x.
\]
For $k\ne j$, column orthogonality now reads
\[
 0=(A^{\mathsf T}A)_{jk}
   =x_jx_k\sum_{\lambda_i=0}t_i^2.
\]
As $x_j\ne0$ and the sum is positive, $x_k=0$ for every $k\ne j$.
Thus all zero-weight rows lie on the $j$th coordinate axis.

Let $q$ be their number. Since each row has squared norm at most $1$
and column $j$ has energy $m+\delta$, we have $q\ge m+1$.
The remaining two columns would then satisfy
\[
 2(m+\delta)
 =\sum_{\lambda_i>0}\sum_{k\ne j}a_{ik}^2
 \le n-q\le2m.
\]
\end{proof}

For $n=3m+2$, the kernel cases need a separate argument
(see \cref{lem:r2-all-active}).  We use positivity of the row weights
in the support bound below.

\subsection{Second-order reduction of row support}
The next arguments apply when a full-support row occurs.
Positive definiteness of each multiplier matrix makes a strict
second-order descent possible whenever too many entries are nonzero.

\begin{lemma}[A full row forces positive definiteness]\label{lem:full-pd}
Suppose every row multiplier is positive in every normalized system.
If $A$ has a row with three nonzero entries, every multiplier matrix $S$
is positive definite.
\end{lemma}
\begin{proof}
Change column signs so that a full row is $u>0$. Since it is saturated,
\[
 \mathbf1^{\mathsf T}u=1,\qquad Su=\lambda\mathbf1,
 \qquad \lambda>0.
\]
Put $M=S-\lambda\mathbf1\mathbf1^{\mathsf T}$. By
\eqref{eq:multiplier-bounds} and $Mu=0$,
\[
 M_{ii}\ge0,\qquad M_{ij}\le\min(M_{ii},M_{jj}).
\]
Indeed, $\lambda=\sum_jS_{ij}u_j\le S_{ii}\sum_j u_j=S_{ii}$.
We first show $M\succeq0$.

If all off-diagonal entries of $M$ are nonpositive, then
\[
 \diag(u)M\diag(u)
 =\sum_{i<j}(-u_i u_j M_{ij})
   (e_i-e_j)(e_i-e_j)^{\mathsf T}\succeq0.
\]
Here we used $Mu=0$. Otherwise choose $i\ne j$ with $M_{ij}>0$.
Since $M_{ii},M_{jj}\ge M_{ij}$, the principal block
\[
 \begin{pmatrix}M_{ii}&M_{ij}\\M_{ij}&M_{jj}\end{pmatrix}
 \succeq0.
\]
For any $z\in\R^3$, subtract $(z_\ell/u_\ell)u$, where
$\{i,j,\ell\}=\{1,2,3\}$. The new vector has $\ell$th coordinate zero,
and its quadratic form under $M$ is unchanged. Thus $M\succeq0$ in
both cases, and $S=M+\lambda\mathbf1\mathbf1^{\mathsf T}\succeq0$.

Suppose $S$ is singular, and choose $x\ne0$ with $Sx=0$. Then
\[
 0=x^{\mathsf T}Sx=x^{\mathsf T}Mx
    +\lambda(\mathbf1^{\mathsf T}x)^2,
 \qquad \mathbf1^{\mathsf T}x=0.
\]
Choose opposite-sign coordinates $x_j,x_k$, one of which is the unique
coordinate of its sign. For the third coordinate $x_\ell$,
$|x_j|+|x_k|>|x_\ell|$. Since all $\lambda_i>0$, we have
$g_i^{\mathsf T}x=0$ for every row. If $a_{ij}a_{ik}<0$, then
\[
 |g_{ij}x_j+g_{ik}x_k|=|x_j|+|x_k|
 >|x_\ell|\ge|g_{i\ell}x_\ell|,
\]
contrary to $g_i^{\mathsf T}x=0$. Hence every
$a_{ij}a_{ik}\ge0$. The full positive row makes one product strict,
so
\[
 0=(A^{\mathsf T}A)_{jk}=\sum_i a_{ij}a_{ik}>0,
\]
contradicting column orthogonality. Therefore $S$ is positive definite.
\end{proof}

\begin{lemma}[Fixed-Gram curves with prescribed derivatives]
\label{lem:fixed-gram-jets}
Let $A\in\R^{n\times3}$ be feasible at level $s>0$. Suppose that
$H,W\in\R^{n\times3}$ satisfy
\[
 A^{\mathsf T}H+H^{\mathsf T}A=0,\qquad
 A^{\mathsf T}W+W^{\mathsf T}A=-2H^{\mathsf T}H.
\]
There is a differentiable curve $\Gamma$ in $\mathcal M_s$, defined on
a neighborhood of $0$, whose derivative is differentiable at $0$ and
for which
\[
 \Gamma(0)=A,\qquad \Gamma'(0)=H,\qquad \Gamma''(0)=W.
\]
\end{lemma}
\begin{proof}
For a tangent matrix $V$, define its skew-symmetric lift by
\[
 L_A(V)=s^{-1}(VA^{\mathsf T}-AV^{\mathsf T})
       -s^{-2}A(A^{\mathsf T}V)A^{\mathsf T}.
\]
The identities $A^{\mathsf T}A=sI_3$ and
$V^{\mathsf T}A=-A^{\mathsf T}V$ give
$L_A(V)^{\mathsf T}=-L_A(V)$ and $L_A(V)A=V$.
Put $\Omega=L_A(H)$. Since $\Omega A=H$,
\[
 A^{\mathsf T}\Omega^2A
 = (\Omega^2A)^{\mathsf T}A=-H^{\mathsf T}H.
\]
It follows that $R=W-\Omega^2A$ is tangent. Thus
$\Psi=L_A(R)$ is skew-symmetric and $\Psi A=R$.

Set
\[
 X(t)=\tfrac12t\Omega+\tfrac14t^2\Psi,\qquad
 Q(t)=(I_n+X(t))(I_n-X(t))^{-1},\qquad
 \Gamma(t)=Q(t)A.
\]
The inverse exists near $0$. Because $X(t)$ is skew-symmetric,
the Cayley transform $Q(t)$ is orthogonal, so
$\Gamma(t)^{\mathsf T}\Gamma(t)=sI_3$.
The entries of $Q(t)$ are rational functions without poles near $0$.
Differentiating at $0$ gives
\[
 Q(0)=I_n,\qquad Q'(0)=\Omega,\qquad
 Q''(0)=\Omega^2+\Psi.
\]
Multiplication by $A$ gives the prescribed derivatives, since
$(\Omega^2+\Psi)A=\Omega^2A+R=W$.
\end{proof}

\Needspace{12\baselineskip}
\begin{lemma}[Strict second-order descent]\label{lem:second-order}
Let $H\in T_A\mathcal M_s$ annihilate every active facet.
If
\begin{equation}\label{eq:strict-curvature}
 \tr(SH^{\mathsf T}H)>0
\end{equation}
for every normalized multiplier matrix, then $A$ is not a local minimizer
of $h$ on $\mathcal M_s$.
\end{lemma}
\begin{proof}
Put $W_0=-s^{-1}AH^{\mathsf T}H$. It satisfies
$A^{\mathsf T}W_0+W_0^{\mathsf T}A=-2H^{\mathsf T}H$.
There exists $K\in T_A\mathcal M_s$ for which every active facet is
negative on $W_0+K$. Otherwise the inhomogeneous version of the
alternative used in \cref{lem:multipliers}, a consequence of Motzkin's
transposition theorem \cite[Ch.~7]{Schrijver1986}, would give
a convex combination of active facets vanishing on the tangent space
and nonnegative on $W_0$. Its multiplier matrix $S$ would satisfy
\[
 \langle AS,W_0\rangle=-\tr(SH^{\mathsf T}H)<0,
\]
a contradiction. Put $W=W_0+K$ and apply
\cref{lem:fixed-gram-jets}. For the resulting curve $\Gamma$ and every
active facet $\ell$,
\[
 \ell(\Gamma(t))
 =1+\tfrac12t^2\ell(W)+o(t^2)<1
\]
for all sufficiently small positive $t$, since $\ell(H)=0$ and
$\ell(W)<0$. Each inactive facet remains strictly below $1$ by
continuity. There are only finitely many facets, so these inequalities
hold simultaneously for all sufficiently small positive $t$. Hence
$h(\Gamma(t))<h(A)$ there, which rules out a local minimum at $A$.
\end{proof}

\begin{proposition}[Support bound in the presence of a full row]
\label{prop:support-bound}
Suppose all row multipliers are positive in every normalized system,
and at least one row of $A$ is full. Let $Q$ count rows with exactly two
nonzero entries and $F$ count rows with exactly three. Then
\begin{equation}\label{eq:support-bound}
 Q+2F\le5.
\end{equation}
\end{proposition}
\begin{proof}
Let $Z$ be the number of zero entries. Restrict tangent vectors $H$ to
vanish at these entries. All active facets of a given row then induce
the same linear equation on $H$. The $n$ row equations have at least
one dependence on the tangent space: for every tangent vector $H$,
\[
 \sum_i\lambda_i g_i^{\mathsf T}h_i
 =\langle AS,H\rangle=0.
\]
Thus their rank is at most $n-1$.  Since
$\dim\mathcal M_s=3n-6$, the space annihilating all active facets
has dimension at least
\[
 3n-6-Z-(n-1)=2n-5-Z.
\]
A nonzero vector $H$ in this space would satisfy
\eqref{eq:strict-curvature} for every multiplier because
all multiplier matrices are positive definite by
\cref{lem:full-pd}. This contradicts \cref{lem:second-order}.
Thus $Z\ge2n-5$. All rows are nonzero. If $P$ is the number of pure
rows, then $n=P+Q+F$ and $Z=2P+Q$, giving \eqref{eq:support-bound}.
\end{proof}

\begin{remark}
We use $Q+2F\le5$ only when a full row exists and all row weights
are positive.  The second-order argument requires the curvature
condition for every normalized multiplier, not just one of them.
\end{remark}

\Needspace{12\baselineskip}
\subsection{Pair rows and the acyclic graph}
Assume all rows are saturated, all multipliers are positive, and there
are no full rows. If a pair row meets columns $j,k$ with magnitudes $a,b$
and product sign $\varepsilon$, then $a+b=1$ and
\begin{equation}\label{eq:pair-stationarity}
 (S_{jj}-\varepsilon S_{jk})a
 =(S_{kk}-\varepsilon S_{jk})b.
\end{equation}
Both factors are strictly positive for any pair that occurs. Indeed,
equality $|S_{jk}|=S_{jj}$ in the sum proving
\eqref{eq:multiplier-bounds} would force a single product sign on every
overlapping row, contrary to column orthogonality and positive row weights.
It follows that $a>b$ precisely when $S_{jj}<S_{kk}$.
The \emph{pair-support graph} has vertex set $\{1,2,3\}$, one vertex
for each column. The pair $\{j,k\}$ occurs as an edge when at least
one row has support exactly $\{j,k\}$. If such rows are unbalanced,
orient the edge $j\to k$ when their larger entries are in column $j$.
The preceding stationarity equation ensures that this orientation is
consistent across rows on the same pair. In the directed graph, a
\emph{source} has no incoming edges and a \emph{sink} has no outgoing
edges. Once balanced rows are excluded, the graph is acyclic because
$S_{jj}<S_{kk}$ along each arrow $j\to k$.

For an edge $j\to k$, let $b_i$ be the magnitude in column $k$ of
row $i$ supported on $\{j,k\}$, so its magnitude in column $j$ is
$1-b_i$. Define the \emph{incoming energy} and \emph{outgoing energy
deficit} of the edge by
\[
 E_{jk}=\sum_{i:\,\supp(a_i)=\{j,k\}} b_i^2,
 \qquad
 D_{jk}=\sum_{i:\,\supp(a_i)=\{j,k\}}(2b_i-b_i^2).
\]
These names refer to squared mass in the smaller coordinates and to
the loss of squared mass from the larger coordinates, respectively,
since $1-(1-b_i)^2=2b_i-b_i^2$.
For a vertex $j$ we write
\[
 E_j^{\mathrm{in}}=\sum_{h\to j}E_{hj},\qquad
 D_j^{\mathrm{out}}=\sum_{j\to h}D_{jh}.
\]
All edge sums concern rows, not geometric lengths in the graph.
\par\smallskip
\begin{figure}[H]
\centering
\includegraphics[width=.67\linewidth]{pair_support_graph.pdf}
\caption{Possible orientation of a pair-support graph on three coordinates.
A displayed arrow $j\to k$ means that pair rows are supported on
$\{j,k\}$ and are larger in coordinate $j$. Some arrows may be absent.}
\label{fig:pair-support-graph}
\end{figure}
\par
The transfer inequality below will be used in both residue classes.
Balanced pair rows are treated separately in the residue analysis.

\Needspace{16\baselineskip}
\subsection{Energy transfer along pair rows}

\begin{lemma}[An edge transfer bound]\label{lem:edge-transfer}
Consider saturated pair rows whose larger coordinate is in one fixed
column and whose smaller coordinate is in another. Suppose
the signed column products sum to zero. Let $E$ be the sum
of the squared smaller coordinates, let $L$ be the sum of
the row losses, and put $D=L+E$. Then
\begin{equation}
 L\ge2E,
 \qquad L\ge2\sqrt{2E}-2E,
 \qquad D\ge3E.
 \label{eq:edge-loss}
\end{equation}
An \emph{edge group} is the set of pair rows supported on one fixed
coordinate pair and oriented in the same direction. The same
inequalities hold when several groups use disjoint sets of rows.
Their coordinate pairs may share a column, and $E,L,D$ are then
summed over all groups. If $D<3/2$, then
\begin{equation}
 E\le B(D):=2\left(1-\sqrt{1-D/2}\right)^2.
 \label{eq:edge-inverse}
\end{equation}
Moreover $0\le B'(D)<1$ for $0\le D<3/2$.
\end{lemma}

\begin{proof}
Write the pair magnitudes as $(1-b_i,b_i)$, where
$0<b_i<1/2$. Column orthogonality makes the positive and negative signed-product
sums equal to the same number $P\ge0$.  Hence
$L=2\sum_i b_i(1-b_i)=4P$.  Since $b_i<1/2$,
we also have $L\ge2E$.  If $P\le1/4$, put
\[
 q=\frac{1-\sqrt{1-4P}}2.
\]
Each product $b_i(1-b_i)$ is at most $P=q(1-q)$.
The function $u(1-u)$ increases on $0\le u\le1/2$, so
$b_i\le q$.  For either sign group,
\[
 \sum b_i^2
 \le\frac{q}{1-q}\sum b_i(1-b_i)
 =\frac{qP}{1-q}=q^2.
\]
Thus $E\le2q^2$, and monotonicity of
$2\sqrt{2E}-2E$ on $0\le E\le1/2$ gives
\[
 L=4q(1-q)\ge2\sqrt{2E}-2E.
\]
If $P>1/4$, then $L>1$, whereas
$2\sqrt{2E}-2E\le1$ for all $E\ge0$.
For edge groups with disjoint sets of rows, add the inequalities
and use $\sum_e\sqrt{E_e}\ge\sqrt{\sum_eE_e}$.

If $D<3/2$, then $E\le D/3<1/2$.  Combining $D=L+E$
with the nonlinear loss estimate yields
$D\ge2\sqrt{2E}-E$.  The right side is strictly increasing for
$0\le E<1/2$.  Solving this inequality on that branch gives
\eqref{eq:edge-inverse}.
Finally
\[
 B'(D)=\frac1{\sqrt{1-D/2}}-1<1
 \qquad(0\le D<3/2).
\]
\end{proof}



\subsection{Normal forms with full-support rows}
The classification below applies in either nonzero residue class once
all row multiplier weights are known to be positive. That hypothesis
is established in \cref{lem:active-r1,lem:r2-all-active}.

\Needspace{16\baselineskip}
\begin{lemma}[Full-row normal forms]
\label{lem:full-row-classification}
Let $A$ be an optimal $n\times3$ feasible matrix with
$n=3m+r$, $r\in\{1,2\}$, and
$A^{\mathsf T}A=(m+\delta)I_3$, where $0<\delta<1/2$.
Assume that every normalized multiplier system for $A$ has strictly
positive row weights and that at least one row is full.  Then
$F\in\{1,2\}$, and after row and column permutations and sign changes,
the non-pure rows have exactly one of the following forms.
\begin{enumerate}
\item One full row and three pair rows:
\begin{equation}
 (a,b,c),\quad(1-x,-x,0),\quad(1-y,0,-y),\quad(0,1-z,-z),
 \label{eq:one-full-form}
\end{equation}
where $a,b,c>0$, $a+b+c=1$, $0<x,y,z<1/2$, and
\[
 x(1-x)=ab,\qquad y(1-y)=ac,\qquad z(1-z)=bc.
\]
In particular, $x<\min(a,b)$, $y<\min(a,c)$, and
$z<\min(b,c)$.
\item Two full rows and one pair row:
\begin{equation}
 (a,b,c),\quad(a,-b,c),\quad(g,0,-h),
 \label{eq:two-full-form}
\end{equation}
where all five parameters are positive,
$a+b+c=g+h=1$, and $2ac=gh$.
\end{enumerate}
\end{lemma}
\begin{proof}
If $F=1$, each of its three nonzero column products must be canceled
by a pair row on the same two columns.  Thus $Q\ge3$, while
$Q+2F\le5$ gives $Q=3$.  Consequently there is exactly one pair row
for each of the three column pairs.  Make the full row positive by
changing column signs.  The product of the two nonzero entries in
each pair row must be negative.  No pair can be balanced, since a
balanced pair has product magnitude $1/4$ whereas
$ab,ac,bc<1/4$.

The stationarity equations \eqref{eq:pair-stationarity} for a pair with positive magnitudes $u,v$
and product sign $\varepsilon$ give
\[
 (S_{jj}-\varepsilon S_{jk})u
   =(S_{kk}-\varepsilon S_{jk})v.
\]
Both factors are nonnegative by the multiplier bounds.  If one
vanished, the equation would force the other to vanish, contrary to
$S\succ0$.  Thus both are positive, and the larger coordinate is the
one with smaller diagonal entry of $S$.  Since the three pairs are
unbalanced, these strict inequalities define an acyclic orientation
of the complete graph on three vertices, hence a transitive ordering.
After relabeling, this gives \eqref{eq:one-full-form}.  For example,
$x(1-x)=ab$ and $a+b<1$ imply $x<\min(a,b)$: writing
$t=\min(a,b)<1/2$, one has $ab<t(1-t)$, and the function
$u\mapsto u(1-u)$ is strictly increasing for $0\le u<1/2$.
The other comparisons are analogous.

If $F=2$, two full rows alone cannot cancel all three column products.
Indeed, the product of the three pairwise coordinate products of any
full row is positive, whereas canceling all three products would
require the corresponding products of the other row all to have
opposite signs.  Thus $Q\ge1$, and $Q+2F\le5$ gives $Q=1$.
Place this pair in columns one and three and make the first full row
$(a,b,c)>0$.  Orthogonality of column two with columns one and three
forces the second full row to be
$(ta,-b/t,tc)$ for some $t>0$.
The two full-row budgets give $a+b+c=1$ and
$t(a+c)+b/t=1$.  Multiplying the latter equality by $t$ yields
\[
 (t-1)\bigl(t(1-b)-b\bigr)=0.
\]
If $t\ne1$, then $t=b/(1-b)$, so the second column of these full rows
contributes $b^2+(1-b)^2\in[1/2,1)$ to its squared norm.
Adding an integer number of unit coordinate rows cannot make this
contribution equal to the fractional part $\delta\in(0,1/2)$.
Thus $t=1$, yielding \eqref{eq:two-full-form}.  Finally,
$F\ge3$ contradicts $Q+2F\le5$.
\end{proof}

% END sections/variational
% BEGIN sections/r1
\section{The residue class \texorpdfstring{$n\equiv1\pmod3$}{n = 1 modulo 3}}
\label{sec:residue-one}

Write $n=3m+1$ and $s=m+\delta$ for an optimal frame.
The constructions give $\delta>0$.  By \cref{lem:dominance,cor:two-dimensional},

\begin{equation}
  0<\delta\le (2+\sqrt3)^{-2}<\frac{72}{1000},\qquad
  E\le\frac{(1-\sqrt{3\delta})^2}{2}<\frac12,
  \label{eq:r1-preliminary}
\end{equation}
Here $E$ is precisely the nondominant squared mass defined in
\cref{lem:dominance}: it sums $a_{ij}^2$ over positions where column
$j$ is not dominant for row $i$. The counts of rows dominant in each
column are $(m+1,m,m)$, up to a permutation of the columns.
By \cref{lem:active-r1}, every row is saturated and has positive weight.

\subsection{The two algebraic parameters}\label{sec:algebraic-parameters}

The constants $\delta_1,\delta_2$ in the main theorem are defined
by two explicit equations below. They arise from small matrices with
two orthogonal columns of the form
\[
 R=\begin{pmatrix}xC&yS\\xS&-yC\end{pmatrix},
 \qquad C^2+S^2=1.
\]
The columns of $R$ have squared norms $x^2$ and $y^2$ and are
orthogonal.  The row-budget equalities for the corresponding frames
lead to two linear systems in $C,S$.  Solving them yields the
following scalar functions on $0\le y\le27/100$:
\begin{align*}
 x_1(y)&=\sqrt{1+y^2},&
 C_1(y)&=x_1(y)(\sqrt2-x_1(y))-y,\\
 &&S_1(y)&=x_1(y)-y(\sqrt2-x_1(y)),\\
 x_2(y)&=\sqrt{2+y^2},&
 C_2(y)&=\tfrac12\bigl(x_2(y)(\sqrt2-y)-y\bigr),\\
 &&S_2(y)&=\tfrac12\bigl(x_2(y)-y(\sqrt2-y)\bigr).
\end{align*}
We use \emph{root isolation} to mean proving that an equation has
exactly one real root in a specified rational interval.
The defining equations for our two exceptional parameters are
\begin{equation}\label{eq:root-equations}
 C_j(y)^2+S_j(y)^2=1\qquad(j=1,2).
\end{equation}
The next lemma proves that each equation has exactly one root
$y_j$ in $[0,27/100]$.  We set $\delta_1=y_1^2$ and $\delta_2=y_2^2$.
The parametrizations are recovered from the row budgets in
\cref{prop:r1-two-full}.

\begin{lemma}[Existence, uniqueness, and isolation]
\label{lem:root-isolation}
For $j=1,2$, equation \eqref{eq:root-equations} has a unique solution
$y_j$ in $[0,27/100]$.  These solutions satisfy
\[
 0<y_1<\frac{17}{100},\qquad
 \frac{347}{2000}<y_2<\frac7{40}.
\]
With $\delta_j=y_j^2$ and $\delta_*=15-4\sqrt{14}$, one has
\begin{equation}
 0<\delta_1<\frac{3}{100}<\delta_2<\frac{49}{1600}<\delta_*.
 \label{eq:exceptional-order}
\end{equation}
\end{lemma}
\begin{proof}
All four functions $C_j,S_j$ are positive on $[0,27/100]$.
For the first pair, this follows, for example, from
$\sqrt2>7/5$, $1\le x_1<26/25$, and $y\le27/100$.
Their derivatives are
\[
 C_1'=\frac{\sqrt2\,y}{x_1}-2y-1,
 \qquad
 S_1'=x_1+\frac{y(1+y)}{x_1}-\sqrt2.
\]
Both are negative: for the second derivative expression use
$x_1<26/25$, $y(1+y)/x_1\le3429/10000$, and $\sqrt2>7/5$.
For the second pair,
\[
 C_2'=\frac12\left(\frac{y(\sqrt2-y)}{x_2}-x_2-1\right),
 \qquad
 S_2'=\frac12\left(\frac{y}{x_2}-\sqrt2+2y\right),
\]
which are negative using $x_2\ge\sqrt2>7/5$ and $y\le27/100$.
Thus $C_j^2+S_j^2$ is strictly decreasing.

At $y=0$, the two squared sums are $4-2\sqrt2>1$ and $3/2>1$.
At $y=17/100$, the first sum is less than one. The second sum is
greater than one at $y=347/2000$ and less than one at $y=7/40$.
These rational endpoint checks are given in
\cref{app:root-arithmetic}. Continuity and strict monotonicity
then give the stated isolations.

Squaring the root bounds gives the first two comparisons in
\eqref{eq:exceptional-order}. The last follows from
\[
 23951^2-14\cdot6400^2=210401>0,
\]
which gives $\sqrt{14}<23951/6400$ and hence
$49/1600<15-4\sqrt{14}$.
\end{proof}

\paragraph{The known four-dimensional value.}
Our parameter $x_0=1+\delta_1$ satisfies
\[
 4x_0^4-28x_0^3-7x_0^2+16x_0+16=0,
 \qquad 1<x_0<\frac43,
\]
which is the polynomial appearing in \cite{Shultz2013}.
For completeness, eliminating the square roots from the defining
equation for $y_1$ gives the necessary relation
\[
 (2y^4+4y^3+y^2-1)(2y^4+4y^3+y^2+8y-1)=0.
\]
The first factor is negative throughout $0<y<17/100$, the
interval containing $y_1$. Hence the second factor vanishes
at $y_1$. Substituting $x_0=1+y_1^2$
then gives the displayed quartic. Equivalently, $\delta_1$ is a root of
$p_1(d)=4(1+d)^4-28(1+d)^3-7(1+d)^2+16(1+d)+16$; this quartic is
irreducible over $\mathbb Q$, and its only other real root is
approximately $6.156$.

Similarly, eliminating $x_2$ and $\sqrt2$ from the defining equation
for $y_2$ (two resultants) gives $p_2(y^2)=0$, with $p_2$ as in
\cref{sec:main}. The polynomial $p_2$ is irreducible over $\mathbb Q$,
and its real roots are approximately
$0.0303$, $1.625$, $2.945$ and $8.190$; hence $\delta_2$ is its
smallest positive root.


\begin{proposition}[Exceptional frames]
\label{prop:exceptional-frames}
There are feasible four-row and seven-row frames with Gram matrices
$(1+\delta_1)I_3$ and $(2+\delta_2)I_3$, respectively.
\end{proposition}
\begin{proof}
In the four-row construction put $y=y_1$, $x=\sqrt{1+y^2}$,
$C=C_1(y)$, and $S=S_1(y)$.  The rows are
\begin{equation}
 \left(\frac{xC}{\sqrt2},\frac{x}{\sqrt2},\frac{yS}{\sqrt2}\right),\quad
 \left(\frac{xC}{\sqrt2},-\frac{x}{\sqrt2},\frac{yS}{\sqrt2}\right),\quad
 (xS,0,-yC),\quad (0,0,1).
 \label{eq:four-frame}
\end{equation}
The equations $C^2+S^2=1$,
$xC+yS=\sqrt2-x$, and $xS+yC=1$ show that all row budgets equal one.
For example, the first two rows both have $\ell^1$-norm
$(xC+x+yS)/\sqrt2=1$, since the negative sign in the second row
disappears under absolute values.  The third row has norm
$xS+yC=1$.  The first and third columns are orthogonal because their two
contributions $xyCS$ and $-xySC$ cancel.  The signs in the first two
rows cancel the other off-diagonal entries.  The diagonal entries are
$x^2,x^2,1+y^2$, as required.

For seven rows put $y=y_2$, $x=\sqrt{2+y^2}$,
$C=C_2(y)$, and $S=S_2(y)$.  Take
\begin{equation}
 \left(\frac{xC}{\sqrt2},\frac{y}{\sqrt2},\frac{yS}{\sqrt2}\right),\quad
 \left(\frac{xC}{\sqrt2},-\frac{y}{\sqrt2},\frac{yS}{\sqrt2}\right),\quad
 (xS,0,-yC),\quad e_2,e_2,e_3,e_3.
 \label{eq:seven-frame}
\end{equation}
Here $xC+yS=\sqrt2-y$ and $xS+yC=1$.  The same calculation gives row
budgets one and Gram matrix $(2+y^2)I_3$.
\end{proof}

\subsection{Full-row cases in the first residue class}
By \cref{lem:active-r1}, the positivity hypothesis of
\cref{lem:full-row-classification} holds for an optimum in this
residue class. We now bound its two possible full-row configurations.

\begin{proposition}[One full row in the first residue class]
\label{prop:r1-one-full}
A $(3m+1)$-row optimum cannot have exactly one full row if $m=1$,
or if $\delta\ge3/100$.
\end{proposition}
\begin{proof}
Let $p_j$ count the pure rows in coordinate $j$.  In the normal form \eqref{eq:one-full-form}, the cancellation
relations and $x,y,z<1/2$ give $y^2<ac$ and $z^2<bc$.
Thus the mixed contribution to column three satisfies
\[
 c^2+y^2+z^2<c^2+c(a+b)=c<1.
\]
Since the total energy in this column is $m+\delta$ with
$0<\delta<1/2$, its pure-row count is $p_3=m$.
The mixed squared mass in the second column is
$b^2+x^2+(1-z)^2$. Since $x<a$ and $a+b<1$, we have
$b^2+x^2<b^2+a^2<(a+b)^2<1$. Also $(1-z)^2<1$.
Hence the second mixed contribution is positive and less than $2$, giving
$p_2\in\{m-1,m\}$. The mixed squared mass in the first column is
$a^2+(1-x)^2+(1-y)^2$, which is positive and less than $3$.
Therefore $m-2\le p_1\le m$.  Since $p_1+p_2+p_3=3m-3$, necessarily
\[
 (p_1,p_2,p_3)=(m-2,m-1,m).
\]
This is impossible for $m=1$.  Otherwise the diagonal equations become
\begin{align}
 \delta&=3a^2-2a-x^2-y^2,\label{eq:r1-full-a}\\
 \delta&=b^2+x^2-2bc-z^2,\label{eq:r1-full-b}\\
 \delta&=c^2+y^2+z^2.\label{eq:r1-full-c}
\end{align}
The first equation gives $a>2/3$ and thus $b+c<1/3$.
Since $z<\min(b,c)<1/6$,
\[
 \frac{z}{c}=\frac{b}{1-z}<\frac25.
\]
If $\delta\ge3/100$, \eqref{eq:r1-full-c} gives
$\delta<(54/25)c^2$ and hence $c>117/1000$.
Equation~\eqref{eq:r1-full-b} gives
$\delta<2b(b-c)$, whence $b>193/1000$.
Consequently $a<69/100$ and $b+c>31/100$.
Moreover, $x>ab>(2/3)b$ and $y>ac>(2/3)c$.  Returning to
\eqref{eq:r1-full-a}, we obtain
\[
 \delta<3\left(\frac{69}{100}\right)^2
           -2\left(\frac{69}{100}\right)
           -\frac29\left(\frac{31}{100}\right)^2
       =\frac{97}{3600}<\frac3{100},
\]
a contradiction.
\end{proof}

\begin{proposition}[Two full rows in the first residue class]
\label{prop:r1-two-full}
In the normal form \eqref{eq:two-full-form}, a $(3m+1)$-row optimum has
$\delta=\delta_1$ or $\delta=\delta_2$.  The latter case requires $m\ge2$.
\end{proposition}
\begin{proof}
Let $p_j$ count the pure rows in column $j$. Every mixed column
has positive squared energy, so $p_j\le m$. The second column has
mixed energy $2b^2\in(0,2)$, giving $p_2=m-1$ or $p_2=m$.
There are three mixed rows, so $p_1+p_2+p_3=3m-2$.
These bounds determine the alternatives below.
\Needspace{12\baselineskip}
The two-by-two matrix
\[
 R=\begin{pmatrix}\sqrt2\,a&\sqrt2\,c\\g&-h\end{pmatrix}
\]
has orthogonal columns. The original two full rows have the common
budget $a+b+c=1$, while the pair row has budget $g+h=1$.
After compression, the two row norms of $R$ are therefore
\[
 \norm{R_{1,*}}_1=\sqrt2(a+c)=\sqrt2(1-b),\qquad
 \norm{R_{2,*}}_1=g+h=1.
\]
In particular, the first row of $R$ need not have unit $\ell^1$-norm.
These are the budget identities used in both cases below.

If $p_2=m-1$, then, after interchanging columns one and three if
necessary, $(p_1,p_3)=(m-1,m)$.  Put $y=\sqrt\delta$ and
$x=\sqrt{1+\delta}$.  Orthogonality and positivity give
\[
 R=\begin{pmatrix}xC&yS\\xS&-yC\end{pmatrix},\qquad C,S>0,\quad C^2+S^2=1.
\]
Since $b=x/\sqrt2$, the row budgets say
$xC+yS=\sqrt2-x$ and $xS+yC=1$.  Solving gives
$C=C_1(y)$ and $S=S_1(y)$.  By \eqref{eq:r1-preliminary},
$y<27/100$, so Lemma~\ref{lem:root-isolation} gives $\delta=\delta_1$.

If $p_2=m$, then $(p_1,p_3)$ is $(m-1,m-1)$ or $(m-2,m)$ up to
interchange.  In the first case $R^{\mathsf T}R=(1+\delta)I_2$,
so $g^2+h^2=1+\delta>1$, contrary to $g+h=1$.
In the second case put $x=\sqrt{2+\delta}$, $y=\sqrt\delta$.
The same parametrization yields
$xC+yS=\sqrt2-y$ and $xS+yC=1$.
Now $C=C_2(y)$, $S=S_2(y)$, and $\delta=\delta_2$.
The pure count $p_1=m-2$ requires $m\ge2$.
\end{proof}

\Needspace{14\baselineskip}
\subsection{Frames with no full row}

\begin{proposition}[The first residue class without full rows]
\label{prop:r1-no-full}
If a $(3m+1)$-row optimum has no full row, then $\delta\le\delta_*$.
Such an optimum is impossible when $m=1$, and when $m=2$ it satisfies
$\delta<3/100$.
\end{proposition}
\begin{proof}
Orient every pair from its dominant coordinate to its nondominant
coordinate.  The multiplier equations make this a directed acyclic
graph.  Let $E_j$ be the total incoming nondominant squared energy at
column $j$.  The two columns with $m$ dominant rows satisfy
$E_j\ge\delta$, because their dominant entries contribute squared
energy at most $m$ whereas their total is $m+\delta$.
In the absence of full rows, column orthogonality holds separately
for each pair of columns, so the edge transfer bound applies.

For an edge with incoming energy $E_e$ and loss $L_e$, Lemma~\ref{lem:edge-transfer} gives
\[
 L_e\ge\phi(E_e),\qquad
 \phi(t)=2\sqrt{2t}-2t.
\]
The function $\phi$ is increasing on $[0,1/2]$ and satisfies
$\phi(u)+\phi(v)\ge\phi(u+v)$.  Since the total nondominant energy is
less than $1/2$ by \eqref{eq:r1-preliminary}, grouping edges by target
and using the two columns with only $m$ dominant rows gives
\[
 1-3\delta=\sum_e L_e
       \ge2\phi(\delta)=4\sqrt{2\delta}-4\delta.
\]
On the preliminary interval for $\delta$, this is equivalent to
$\delta\le15-4\sqrt{14}$.

Every occurring edge needs at least two pair rows with opposite
product signs, since its signed products must cancel.  Every source
of the acyclic graph must be the unique column with $m+1$ dominant
rows: a source with only $m$ dominant rows has no incoming energy,
so its squared column norm would be at most $m$.  Both other columns
need incoming energy, and hence each must be reached from that source.

When $m=1$, the source has only two dominant rows.  It cannot branch
into two distinct edges, since that would require at least four rows.
If the two columns with only $m$ dominant rows instead form a chain,
the middle vertex
would need two outgoing dominant rows but has only one.  Both
possibilities are impossible.

When $m=2$, the source has three dominant rows, so it still cannot
branch.  Both deficient vertices must receive energy, hence the
orientation is a chain $1\to2\to3$.  The middle vertex has exactly
two dominant rows, and both belong to its outgoing edge.  Write their
nondominant magnitudes as $b_1,b_2$.  The sink has two pure rows,
so $b_1^2+b_2^2=\delta$.  The middle column has energy
\[
 2+\delta=(1-b_1)^2+(1-b_2)^2+E_{12}.
\]
Using $b_1^2+b_2^2=\delta$ gives
\[
 E_{12}=2(b_1+b_2)\ge2\sqrt\delta.
\]
If $\delta\ge3/100$, this would imply
$E\ge2\sqrt{3/100}>3/10$, whereas
\eqref{eq:r1-preliminary} gives
$E\le(1-3/10)^2/2=49/200<3/10$.
\end{proof}

\Needspace{12\baselineskip}
\begin{theorem}[Exact values in the first residue class]
\label{thm:r1}
For every integer $m\ge1$, the exact values in the first nonzero
residue class are
\[
 F_3(4)=1+\delta_1,\qquad F_3(7)=2+\delta_2,
 \qquad F_3(3m+1)=m+15-4\sqrt{14}\quad(m\ge3).
\]
\end{theorem}
\begin{proof}
The frames in Proposition~\ref{prop:exceptional-frames} give the first
two lower bounds.  The four-pair star construction gives the third.
Consider an optimum and apply Lemma~\ref{lem:active-r1}.
If it has a full row, the support bound and
Lemma~\ref{lem:full-row-classification} leave the one-full and two-full
forms.  Propositions~\ref{prop:r1-one-full} and
\ref{prop:r1-two-full}, together with
\eqref{eq:exceptional-order}, show that the only possible maximizing
full-row values are $\delta_1$ at $m=1$ and $\delta_2$ at $m=2$.
For $m\ge3$ they are strictly below the available star value.
If there is no full row, Proposition~\ref{prop:r1-no-full} excludes
$m=1$, falls below the available seven-row value at $m=2$, and gives
the sharp upper bound $\delta\le\delta_*$ for $m\ge3$.
\end{proof}

% END sections/r1
% BEGIN sections/r2
\section{The residue class \texorpdfstring{$n\equiv2\pmod3$}{n = 2 mod 3}}
\label{sec:residue-two}

Write
\[
 \Delta=10-4\sqrt6=(\sqrt6-2)^2.
\]
Let $m\ge1$, $n=3m+2$, and let $A$ be a maximizing frame:

\[
 A^{\mathsf T}A=sI_3,\qquad s=m+\delta,\qquad
 \sum_{j=1}^3|a_{ij}|\le1.
\]
Use $AS=G$ and $G_i=\lambda_i g_i$ as in \cref{lem:multipliers}.
Here $\lambda_i\ge0$, $\sum_i\lambda_i=1$, and
$g_i\in\partial\|a_i\|_1$.  We call row $i$ a
\emph{zero-weight row} when $\lambda_i=0$ in the chosen multiplier
system. This describes a multiplier weight, not a zero vector.
Such a row need not be zero and need not saturate its row budget. A \emph{large entry} is an entry with $|a_{ij}|>1/2$.
This condition is different from dominance in \cref{lem:dominance},
which compares one entry with the sum of the remaining magnitudes.
For nonnegative numbers $r_1,r_2,r_3$, the \emph{triangle inequalities}
are $r_j\le r_k+r_\ell$ for all distinct $j,k,\ell$.
They are \emph{strict} if all three inequalities are strict and
\emph{degenerate} if at least one is an equality. We apply these
terms to the absolute coordinate values of a kernel vector.

\subsection{An extremal construction and preliminary bounds}

The five-row matrix $A_5$ in \eqref{eq:five-core} has row
$\ell^1$-norms equal to one and satisfies
$A_5^{\mathsf T}A_5=(1+\Delta)I_3$.
Appending $m-1$ copies of $I_3$ gives
\begin{equation}
 \F(3m+2)\ge m+\Delta.
 \label{eq:r2-construction}
\end{equation}

\begin{proposition}\label{prop:r2-column-bounds}
For any maximizing frame at level $m+\delta$ in dimension $3m+2$,
\[
 \frac15<\Delta\le\delta\le\frac49.
\]
Every row is nonzero.  Each column has at least $m$ entries whose
absolute values exceed $1/2$.
\end{proposition}

\begin{proof}
The trace bound gives $\delta\le2/3<1$.  For numbers
$0\le c_1,\ldots,c_N\le1$ with
$\sum_i c_i^2=m+\delta$, their sum is at least
$m+\sqrt\delta$.  To see this, put $t_i=c_i^2$.  Minimizing
$\sum_i\sqrt{t_i}$ over $0\le t_i\le1$ with fixed
$\sum_i t_i=m+\delta$ reduces, by concavity of the square-root
function, to a vector with $m$ entries equal to $1$, one equal to
$\delta$, and the rest zero.  Its square-root sum is
$m+\sqrt\delta$.
Apply this observation to each column and add the row bounds:
\[
 3m+3\sqrt\delta\le\sum_{i,j}|a_{ij}|\le3m+2.
\]
Thus $\delta\le4/9$. A zero row would instead give
$3\sqrt\delta\le1$, contrary to \eqref{eq:r2-construction}.
The inequality $\Delta>1/5$ follows from $6<(49/20)^2$.

Let $L_i=1-\|a_i\|_2^2$ and $L=\sum_iL_i=2-3\delta$.
If $c=|a_{ij}|\le1/2$, then
\[
 L_i\ge1-c^2-(1-c)^2=2c(1-c)\ge2c^2.
\]
Consequently, if a column had at most $m-1$ entries larger than
$1/2$, its squared norm would be at most
\[
 m-1+\frac L2=m-\frac32\delta<m+\delta,
\]
a contradiction. A row has at most one entry larger than $1/2$.
\end{proof}

\subsection{Positivity of every row weight}

The obstacle in the second residue class is the possible vanishing
of a row multiplier. We reduce this to the support of a single
zero-weight row. The two-coordinate case is settled by a column
count, while a full-support row leads to strict or degenerate
triangle inequalities in the kernel.

\begin{lemma}[Kernel reduction for zero-weight rows]
\label{lem:r2-kernel-reduction}
Let $m\ge1$, let $A$ be a maximizing feasible frame with $3m+2$
rows, and let $(\lambda,g,S)$ be a normalized multiplier system. If
$\lambda_{i_0}=0$ for some row index $i_0$, write $r=a_{i_0}$.
Then $r\ne0$, $S$ is positive semidefinite of rank two, and
$\ker S=\R r$. The row $r$ has no strictly dominant coordinate.
Moreover, $i_0$ is the unique zero-weight index. Every other row
is saturated and satisfies $g_i\cdot r=0$.
\end{lemma}
\begin{proof}
Fix a normalized multiplier with $\lambda_{i_0}=0$, and write
$r=a_{i_0}$. By \cref{prop:r2-column-bounds},
$s=m+\delta$ with $1/5<\delta\le4/9$. Its matrix $S$ satisfies
\[
 S_{jj}\ge0,\qquad |S_{jk}|\le\min(S_{jj},S_{kk}),
 \qquad \operatorname{tr}S=1/s.
\]
The rank-one exclusion in \cref{lem:rank-one} gives
$\operatorname{rank}S\ge2$. The assumed zero-weight row $r$
is nonzero by \cref{prop:r2-column-bounds}. Since $AS=G$ and
$\lambda_{i_0}=0$, we have $rS=0$, so $S$ is singular.
Consequently $\operatorname{rank}S=2$.  Its principal minors of
orders one and two are nonnegative, hence $S$ is positive
semidefinite.  All zero-weight rows lie on the one-dimensional kernel
$\R x=\ker S$. Since $r\ne0$ is in this kernel,
$\ker S=\R r$. Every positive-weight row satisfies
$g_i\cdot r=0$.

Suppose $x$ has a strictly dominant coordinate $j$, so that
$|x_j|>\sum_{k\ne j}|x_k|$.
Then $a_{ij}=0$ in every
positive-weight row. Orthogonality of columns forces all other
coordinates of $x$ to vanish. At least $m+1$ zero-weight rows are needed to supply the
squared column norm $s>m$, since each row has squared norm at most
$1$.  The remaining at most $2m+1$ rows use only the other two
columns.  Appending zero rows if necessary produces a two-dimensional
feasible frame with $2m+1$ rows and level $s=m+\delta$.
The two-dimensional residue estimate (Corollary~\ref{cor:two-dimensional})
\[
 F_2(2m+1)\le m+\frac18
\]
contradicts $\delta>1/5$.

It follows that $x$ has no strictly dominant coordinate.
A vector $v$ parallel to $x$ satisfies
\[
 \|v\|_2^2\le2\sum_{j<k}|v_jv_k|.
\]
Indeed, if $a$ is its largest coordinate magnitude and $b$ is the
sum of the other magnitudes, then $a\le b$ and
$\|v\|_2^2-2\sum_{j<k}|v_jv_k|\le(a+b)(a-b)\le0$.
Let $q$ be the number of zero-weight rows, and let $E_0$ and
$P_0$ denote their total squared norm and total absolute pair
products. Their loss is $q-E_0$. Their pair products have a
fixed sign for each column pair, so the positive-weight rows
must have total absolute pair products at least $P_0$.
Their loss is at least $2P_0$. Thus
\[
 L\ge q-E_0+2P_0\ge q.
\]
Since $L=2-3\delta<7/5$ and at least one such row exists,
there is exactly one zero-weight row, namely $r$.
All other rows have positive weight and are saturated.
\end{proof}

\Needspace{12\baselineskip}
\begin{lemma}[Excluding two-coordinate zero-weight rows]
\label{lem:r2-zero-weight-pair}
For $m\ge1$, no normalized multiplier system of a maximizing frame
with $3m+2$ rows has a zero-weight row supported on exactly two
coordinates.
\end{lemma}
\begin{proof}
Suppose a zero-weight row $r$ has support of size two.
By \cref{lem:r2-kernel-reduction}, it is the only zero-weight row.
After changing column signs and permuting columns, write
$r=(t,t,0)$, where $0<t\le1/2$.
In each positive-weight row, the first two coordinates have
opposite signs whenever both are nonzero. This follows from
$g_i\cdot r=0$ and $r=(t,t,0)$. Orthogonality consequently gives
\[
 \sum_{i\ne r}|a_{i1}a_{i2}|=t^2.
\]
All these rows are saturated. Define
\[
 P_3=\sum_{i\ne r}|a_{i3}|(|a_{i1}|+|a_{i2}|).
\]
The exact loss identity is
\begin{equation}
 L=1+2P_3,\qquad 2P_3=1-3\delta.
 \label{eq:r2-two-kernel-loss}
\end{equation}
Let $p_j$ count the entries larger than $1/2$ in column $j$.
The row $r$ contributes to none of these counts. If $p_3=m$,
the small entries in the third column carry squared mass at
least $\delta$. In a saturated row, a small third coordinate
$c$ contributes $c(1-c)\ge c^2$ to $P_3$.
Hence $P_3\ge\delta$, and \eqref{eq:r2-two-kernel-loss}
would give $\delta\le1/5$.

Hence $p_3\ge m+1$.  By \cref{prop:r2-column-bounds},
$p_1,p_2\ge m$.  The zero-weight row $r$ has no entry larger
than $1/2$, and the remaining $3m+1$ rows have at most one each.
Thus the only remaining count is
$(p_1,p_2,p_3)=(m,m,m+1)$, and every row other than $r$ has
exactly one large coordinate.

Let $J$ be the rows whose large coordinate is the third.
For $i\ne r$ put $q_i=|a_{i1}|+|a_{i2}|$.  Since the signed first-two
column products of the positive-weight rows sum to $-t^2$,
\[
 \begin{aligned}
 2s
 &=2t^2+\sum_{i\ne r}(a_{i1}^2+a_{i2}^2)\\
 &=\sum_{i\ne r}q_i^2.
 \end{aligned}
\]
For the $2m$ rows dominated by the first or second coordinate,
$q_i\le\|a_i\|_1=1$.  For $i\in J$, saturation and a large third
entry give $q_i<|a_{i3}|$, so $q_i^2\le q_i|a_{i3}|$.
Therefore
\[
 2s\le2m+\sum_{i\in J}q_i^2\le2m+P_3.
\]
Thus $2\delta\le P_3$, which together with
\eqref{eq:r2-two-kernel-loss} gives $\delta\le1/7$,
contrary to $\delta>1/5$.
\end{proof}

\Needspace{15\baselineskip}
\begin{lemma}[Pair-support structure around a zero-weight full row]
\label{lem:r2-full-kernel-structure}
Let $m\ge1$. Suppose that a normalized multiplier system of a
maximizing frame with $3m+2$ rows has a zero-weight row $r$
with full support.
After column sign changes, $r_j>0$ and $r_j\le r_k+r_\ell$
for each permutation of $\{1,2,3\}$.
All other rows are saturated and supported on at most two columns,
and every coordinate pair occurs. The oriented pair-support graph
is an acyclic tournament. Every pair row has a unique coordinate
larger than $1/2$ in absolute value. The large-entry counts are
$(m+1,m,m)$, up to permutation, and the source has $m+1$ such entries.
\end{lemma}
\begin{proof}
Suppose $r$ is a zero-weight row with all three coordinates
nonzero. By \cref{lem:r2-kernel-reduction}, it is the only
zero-weight row and every other row is saturated.
Change column signs so that $r_j>0$ for all $j$.
Its coordinates satisfy the triangle inequalities, and
$r_j\le1/2$, since $\rho:=r_1+r_2+r_3\le1$.
We first exclude a second full-support row $u$.
Indeed, if $\sigma=\operatorname{sign}u$, then
$\sigma\cdot r=0$ because $g_u\cdot r=0$ and $u$ is full.
Consider the row rotation
\[
 u(\theta)=u\cos\theta+r\sin\theta,\qquad
 r(\theta)=r\cos\theta-u\sin\theta.
\]
This rotation preserves the contribution
$u^{\mathsf T}u+r^{\mathsf T}r$ to the Gram matrix.
For sufficiently small $\theta$, $u(\theta)$ has the same coordinate
signs as $u$, and $r(\theta)$ has positive coordinates.  Thus
\[
 \begin{aligned}
 \|u(\theta)\|_1
 &=\cos\theta\,\sigma\cdot u
   +\sin\theta\,\sigma\cdot r=\cos\theta<1,\\
 \|r(\theta)\|_1
 &=\rho\cos\theta-\sin\theta\sum_j u_j.
 \end{aligned}
\]
Choose the sign of a sufficiently small nonzero $\theta$ so that
$\sin\theta\sum_j u_j\ge0$.  Since $\rho\le1$, the second norm
is also strictly less than $1$.  The rows $u,r$ are linearly
independent, because $\sigma\cdot u=1$ and $\sigma\cdot r=0$,
and hence their rotated versions remain independent.
At the new maximizing frame, these two rows are slack, so any
normalized multiplier system gives them weight zero.  Its multiplier
matrix must annihilate both independent rows and thus has rank at most
one, contradicting \cref{lem:rank-one}.

All remaining rows are therefore pure or supported on a pair
of coordinates.  Since the zero-weight row $r$ has all three
coordinates positive, column orthogonality forces at least one pair
row on each of the three column pairs, to cancel $r_jr_k>0$.
Every principal $2\times2$ minor of $S$ is strictly positive.
To see this, recall that a rank-two positive semidefinite
$3\times3$ matrix has
$\operatorname{adj}(S)=c\,rr^{\mathsf T}$ for some $c>0$, because
$\ker S=\R r$.  As every $r_j>0$, all diagonal cofactors, which are
precisely the principal minors of order two, are positive.
For a pair row, the multiplier equation reads
\begin{equation}
 (S_{jj}-\varepsilon S_{jk})a
   =(S_{kk}-\varepsilon S_{jk})b,
 \qquad a+b=1,
 \label{eq:r2-pair-stationarity}
\end{equation}
where $a,b>0$ and $\varepsilon$ is the product of the two signs.
If $S_{jj}=S_{kk}$, all rows on this pair have magnitudes
$(1/2,1/2)$. Their signed products are then integer multiples
of $1/4$, and cannot cancel $0<r_jr_k<1/4$.
Thus the three diagonal entries of $S$ are distinct.
Orient each pair from its larger coordinate to its smaller
coordinate. Equation \eqref{eq:r2-pair-stationarity} orders
these edges by increasing diagonal entries of $S$.
The resulting graph is acyclic. Every positive-weight row has a
unique large coordinate, and their counts are $(m+1,m,m)$.

For each oriented edge, use the energy terms $E_{jk},D_{jk}$ defined
before \cref{fig:pair-support-graph}. The source of this acyclic graph must be the column with
$m+1$ large entries.  Indeed, orthogonality of columns $j,k$
implies the absolute product sum of their pair rows is at least
$r_jr_k$.  Since $2b_i-b_i^2\ge2b_i(1-b_i)$, an outgoing edge from
$j$ to $k$ has deficit $D_{jk}\ge2r_jr_k$.
If a source had only $m$ large entries, it would have no incoming
energy, so its squared norm would be at most
\[
 m+r_j^2-2r_j(r_k+r_\ell)\le m,
\]
contradicting $m+\delta>m$.
\end{proof}

\begin{lemma}[Excluding strict-triangle zero-weight full rows]
\label{lem:r2-strict-triangle}
For $m\ge1$, a maximizing frame with $3m+2$ rows admits no normalized
multiplier system with a zero-weight full-support row $r$ whose
coordinate magnitudes satisfy all three strict triangle inequalities.
\end{lemma}
\begin{proof}
Change column signs so that $r_j>0$ for all $j$.
The pair-support graph is an acyclic tournament by
\cref{lem:r2-full-kernel-structure}. The strict triangle
inequalities imply that every pair row has opposite coordinate signs.  Otherwise the two nonzero entries
of a pair-row subgradient would contribute $r_j+r_k$ in absolute
value to $g_i\cdot r$, which the remaining coordinate could not cancel
because $r_j+r_k>r_\ell$.  Consequently all pair products on
$j,k$ are negative and their absolute sum is exactly $r_jr_k$ by
orthogonality.  Since $b_i^2<b_i(1-b_i)$, it follows that
$E_{jk}\le r_jr_k$ and $D_{jk}\ge2r_jr_k$.
Let $j,k,\ell$ be source, middle, and sink. Their energy
balances imply
\[
 \delta\le r_\ell^2+r_jr_\ell+r_kr_\ell\le r_\ell,
\]
and
\[
 \delta\le r_k^2+r_jr_k-2r_kr_\ell
 \le r_k(1-3r_\ell)
 \le\frac{1-3\delta}{2}.
\]
The intermediate bound $\delta\le r_k(1-3r_\ell)$ and
$\delta>0$, $r_k>0$ imply $1-3r_\ell>0$.
Thus $r_k\le1/2$ and $r_\ell\ge\delta$ justify
the last inequality. This contradicts $\delta>1/5$.
\end{proof}

\begin{lemma}[Excluding degenerate-triangle zero-weight full rows]
\label{lem:r2-degenerate-triangle}
For $m\ge1$, a maximizing frame with $3m+2$ rows admits no normalized
multiplier system with a zero-weight full-support row $r$ whose
coordinate magnitudes satisfy a degenerate triangle equality.
\end{lemma}
\begin{proof}
Change column signs so that $r_j>0$ for all $j$.
Suppose instead that one triangle inequality is an equality.
By \cref{lem:r2-full-kernel-structure}, the orientation and
large-entry counts are as stated there. Relabel so that
$r_1=r_2+r_3\le1/2$. Products on pairs $12$ and $13$ are
negative. If $P$ is the sum of the positive products on pair
$23$, the negative products there have total $P+r_2r_3$.
The exact loss is
\begin{equation}
 L=1+4r_2r_3+4P,
 \qquad r_2r_3+P<\frac1{10}.
 \label{eq:r2-degenerate-kernel-loss}
\end{equation}
For pair $12$, each smaller coordinate $b_i<1/2$ satisfies
\[
 b_i(1-b_i)\le r_1r_2\le r_2(1-r_2),
\]
so $b_i\le r_2$ and
\[
 E_{12}\le\frac{r_2}{1-r_2}r_1r_2\le r_2^2.
\]
Similarly $E_{13}\le r_3^2$. These bounds are independent of
the orientation of the respective pairs.
On pair $23$, \eqref{eq:r2-degenerate-kernel-loss} gives
$b_i(1-b_i)<1/10<7/64$, hence $b_i<1/8$. Therefore
\begin{equation}
 E_{23}<\frac17(r_2r_3+2P)<\frac1{35}.
 \label{eq:r2-small-cross-edge}
\end{equation}
Since the directed triangle has exactly one sink, it remains to
consider the three possible locations of that sink.

\emph{Sink $1$.}  The middle coordinate is $j=2$ or $3$.
Its incoming pair energy comes only from the other coordinate among
$2,3$ and is at most $E_{23}<1/35$.  The energy balance for the
middle coordinate gives
\[
 \delta=r_j^2+E_{23}-D_{j1}
 \le r_j^2+E_{23}-2r_jr_1<E_{23}<\frac1{35},
\]
where $r_j\le r_1$ was used.

\emph{Sink $2$.}  Its incoming pair energy is at most
$E_{12}+E_{23}$, hence
\[
 \delta\le r_2^2+E_{12}+E_{23}
 \le2r_2^2+\frac1{35}.
\]
Because $\delta>1/5$, this forces $r_2>1/4$, and therefore
$r_3=r_1-r_2<1/4$.
If coordinate $1$ is the source, then $3$ is the middle coordinate.
Its energy increment satisfies
$\delta\le r_3^2+E_{13}\le2r_3^2<1/8$.
If coordinate $3$ is the source, then $1$ is the middle coordinate,
with
\[
 \delta\le r_1^2+E_{13}-D_{12}
 \le r_1^2+r_3^2-2r_1r_2
 =2r_3^2-r_2^2\le r_3^2<\frac1{16}.
\]

\emph{Sink $3$.}  Interchanging coordinates $2$ and $3$ in the
preceding argument gives the same contradiction.
Each possible sink leads to a contradiction with $\delta>1/5$.
\end{proof}

\begin{lemma}[Positive row weights in residue two]
\label{lem:r2-all-active}
Let $m\ge1$ and let $A$ be any maximizing feasible frame in
dimension $3m+2$. Every normalized multiplier system for $A$
has $\lambda_i>0$ for
all $i$. Consequently all row constraints are saturated.
\end{lemma}
\begin{proof}
Fix an arbitrary normalized multiplier system and suppose that
$\lambda_{i_0}=0$. By \cref{lem:r2-kernel-reduction}, the
nonzero row $r=a_{i_0}$ has no strictly dominant coordinate, so its
support has size two or three. Support of size two contradicts
\cref{lem:r2-zero-weight-pair}. In the full-support case, the
non-strict triangle inequalities follow from kernel reduction.
If all three triangle inequalities are strict, this contradicts
\cref{lem:r2-strict-triangle}. If one is an equality, it contradicts
\cref{lem:r2-degenerate-triangle}. We have ruled out every possible support for a zero-weight row.
Thus every weight is positive, and complementarity in
\cref{lem:multipliers} gives
$\|a_i\|_1=1$ for every row.
\end{proof}

\subsection{Excluding full-support rows}
Now that \cref{lem:r2-all-active} supplies positive row weights,
the common normal forms in \cref{lem:full-row-classification}
can be used to exclude full-support rows in this residue class.

\begin{proposition}[Exclusion of full rows in the second residue class]
\label{prop:r2-full-exclusion}
Let $A$ be a maximizing feasible matrix in dimension $3m+2$,
with $A^{\mathsf T}A=(m+\delta)I_3$ and
$1/5<\delta\le4/9$.  Suppose every normalized multiplier system for
$A$ has strictly positive row weights.  Then $A$ has no full row.
\end{proposition}
\begin{proof}
In the one-full-row form, the third mixed column has squared mass
$c^2+y^2+z^2<1$, as in the proof of
\cref{prop:r1-one-full}, so $p_3=m$. In the second column,
$b^2+x^2+(1-z)^2<2$, giving $p_2\in\{m-1,m\}$.
The first column contains the two pair-row entries $1-x$ and $1-y$.
Since $x,y<1/2$ and $\delta\le4/9$, their squared sum satisfies
\[
 (1-x)^2+(1-y)^2>\frac12>\delta.
\]
Thus $p_1\le m-1$, since $p_1=m$ would leave only $\delta$
units of squared mass for all non-pure rows in that column.
The total pure count $p_1+p_2+p_3=3m-2$ now leaves exactly
$(p_1,p_2)=(m-2,m)$ or $(m-1,m-1)$.
For the first possibility,
$\delta=3a^2-2a-x^2-y^2>1/5$ implies $a>3/4$.
Then $c<1/4$ and
$\delta=c^2+y^2+z^2<3c^2<3/16$, a contradiction.
For the second possibility,
\[
 \delta=1+3a^2-2a-x^2-y^2,
 \quad \delta=b^2+x^2-2bc-z^2,
 \quad \delta=c^2+y^2+z^2.
\]
The last two equations give $c>1/4$ and $b>9/20$, using
$\delta<3c^2$ and $\delta<2b(b-c)$.
Thus $a<3/10$.  Since $x,y<a$, the first equation yields
$\delta>(1-a)^2>49/100>4/9$, again a contradiction.

In the two-full-row form, every mixed column has positive squared
energy, so $p_j\le m$ for all $j$. Also, $p_2\in\{m-1,m\}$.
The pure count is $p_1+p_2+p_3=3m-1$.
If $p_2=m-1$, then $p_1=p_3=m$ and
$R^{\mathsf T}R=\delta I_2$.  This gives
$g^2+h^2=\delta<1/2$, contrary to $g+h=1$.
If $p_2=m$, then $(p_1,p_3)=(m-1,m)$ up to interchange.
With $y=\sqrt\delta$, $x=\sqrt{1+y^2}$, the row equations give
\[
 C=x(\sqrt2-y)-y,\qquad S=x-y(\sqrt2-y),\qquad C^2+S^2=1.
\]
Here $\sqrt{1/5}<y\le2/3$.  The function $C$ is strictly decreasing,
since
\[
 C'=\frac{\sqrt2\,y-2y^2-1}{\sqrt{1+y^2}}-1<0.
\]
At the left endpoint, the bounds
$\sqrt{6/5}<11/10$, $\sqrt2<10/7$, and $\sqrt{1/5}>4/9$
give $C<67/105<16/25$.
The function $S$ is strictly convex because
$S''=2+(1+y^2)^{-3/2}>0$.
At the left endpoint $S<7/10$, while at $y=2/3$,
using $\sqrt{13}<181/50$ and $\sqrt2>7/5$, one gets
$S<323/450<18/25$.
Both $C,S$ are positive for the putative frame.  Hence
\[
 C^2+S^2<\left(\frac{16}{25}\right)^2+
                 \left(\frac{18}{25}\right)^2
          =\frac{116}{125}<1,
\]
the final contradiction.
\end{proof}

\subsection{The upper bound when no row has full support}

\begin{proposition}\label{prop:r2-no-full}
If a maximizing frame in dimension $3m+2$ has no full-support
row, then $\delta\le\Delta$.
\end{proposition}

\begin{proof}
By \cref{lem:r2-all-active}, all row weights are positive and all
rows are saturated.  Suppose, for a contradiction, that some pair row is balanced,
with nonzero magnitudes $(1/2,1/2)$.  By \eqref{eq:pair-stationarity}, the two corresponding
diagonal entries of $S$ are equal.  Its positive stationarity factors
then force every pair row on these columns to be balanced.
Column orthogonality requires at least two such rows with opposite
product signs.  Each has loss $1/2$, so their total loss is at least
$1$.

Neither of these balanced rows has an entry larger than $1/2$.
Of the $3m+2$ rows, at most $3m$ can therefore carry a large entry.
By \cref{prop:r2-column-bounds}, each column must have exactly $m$
large entries.  The third column, outside the balanced pair,
requires at least $\delta$ of squared energy from small entries
in other rows.  Their total loss is at least $2\delta$ by the
small-entry inequality in \cref{prop:r2-column-bounds}.
Thus
\[
 L=2-3\delta\ge1+2\delta,
\]
contrary to $\delta>1/5$.

Thus every row has a unique coordinate larger than $1/2$.
The oriented pair graph is acyclic by
\eqref{eq:pair-stationarity}.  Since every column has at least
$m$ large entries and there are $3m+2$ rows, the counts are either
$(m+2,m,m)$ or $(m+1,m+1,m)$, up to permutation.
For a column with $p_j$ large entries, each pure row contributes
squared mass $1$. A pair row directed away from this column
reduces that contribution by $2b_i-b_i^2$, while an incoming
pair row contributes $b_i^2$. Therefore the column-energy balance is
\begin{equation}
 m+\delta=p_j-D_j^{\mathrm{out}}+E_j^{\mathrm{in}}.
 \label{eq:r2-vertex-balance}
\end{equation}

Suppose first that the counts are $(m+2,m,m)$.
Each of the two columns with count $m$ has incoming energy
at least $\delta$. For each group of edges entering one of those columns, let $E$
and $L$ denote its incoming energy and total row loss.  By \cref{lem:edge-transfer},
\[
 L\ge
 \begin{cases}
  2\sqrt{2E}-2E\ge2\sqrt{2\delta}-2\delta,&
      \delta\le E\le1/2,\\
  2E\ge1\ge2\sqrt{2\delta}-2\delta,&E\ge1/2.
 \end{cases}
\]
The groups entering the two columns with exactly $m$ large entries
are disjoint. Hence
\[
 2-3\delta\ge4\sqrt{2\delta}-4\delta.
\]
Since $\delta\le4/9$, this yields
$\delta\le14-8\sqrt3<1/5$, a contradiction.

The counts must therefore be $(m+1,m+1,m)$.
Choose a topological ordering $1,2,3$ of the pair graph.
The source cannot have only $m$ large entries. Indeed, it has
no incoming energy, while its outgoing deficit is nonnegative,
so \eqref{eq:r2-vertex-balance} would give $m+\delta\le m$. Put $D_0=1-\delta$.
At the source,
\[
 D_{12}+D_{13}=D_0,\qquad
 E_{12}+E_{13}\le B(D_0).
\]
Because $\delta\ge\Delta$ and $B$ is increasing,
\begin{equation}
 B(D_0)\le B(1-\Delta)=\frac\Delta2\le\frac\delta2.
 \label{eq:r2-source-cap}
\end{equation}
The equality follows by substituting
$\Delta=(\sqrt6-2)^2$ into
$B(D)=2(1-\sqrt{1-D/2})^2$.  Thus all outgoing edges of
the source together carry at most $\delta/2$ of incoming energy.
If the middle vertex had count $m$, then
$\delta=E_{12}-D_{23}\le E_{12}\le\delta/2$.
Thus the middle vertex has count $m+1$ and the sink has
count $m$.

The middle and sink balances now give
\[
 D_{23}=D_0+E_{12}<1,
 \qquad \delta=E_{13}+E_{23}.
\]
To verify the first strict inequality,
\eqref{eq:r2-source-cap} gives
$E_{12}\le B(D_0)\le\delta/2$, and therefore
$D_{23}\le1-\delta/2<1$ because $\delta>0$.
Using the transfer bound, we have
\[
 \delta=E_{13}+E_{23}
 \le B(D_0)-E_{12}+B(D_0+E_{12}).
\]
Since $B'(D)<1$ on $0\le D<3/2$ and
$D_0+E_{12}=D_{23}<1$, the mean value theorem implies
$B(D_0+E_{12})-B(D_0)\le E_{12}$.  Consequently
\[
 \delta\le2B(D_0)
 =4\left(1-\sqrt{\frac{1+\delta}{2}}\right)^2.
\]
Both sides of
$\sqrt\delta/2\le1-\sqrt{(1+\delta)/2}$ are nonnegative.
Squaring and rearranging yields
$\delta+4\sqrt\delta-2\le0$, hence
$\sqrt\delta\le\sqrt6-2$ and $\delta\le\Delta$.
\end{proof}

\begin{theorem}\label{thm:r2-sharp}
For every integer $m\ge1$,
\[
 \F(3m+2)=m+10-4\sqrt6.
\]
\end{theorem}

\begin{proof}
The construction \eqref{eq:r2-construction} gives the lower
bound. For a maximizing frame, Lemma~\ref{lem:r2-all-active}
gives positive row weights. If a full-support row is present,
Lemma~\ref{lem:full-pd} and Proposition~\ref{prop:support-bound}
give $Q+2F\le5$. Lemma~\ref{lem:full-row-classification} and
Proposition~\ref{prop:r2-full-exclusion} then rule out that case.
With no full-support row, Proposition~\ref{prop:r2-no-full}
gives the matching upper bound.
\end{proof}

\begin{proof}[Completion of the proof of \cref{thm:sharp-values}]
The divisible case is \cref{prop:matrix-model}. The first-residue
analysis proves the two exceptional values and the stable star value.
The second-residue analysis proves the value supplied by the five-row
core. These exhaust the integers $n\ge3$. The recurrence was deduced
from precisely these branches in \cref{cor:eventual-recurrence}.
\end{proof}

% END sections/r2
% BEGIN sections/consequences
\Needspace{18\baselineskip}
\section{Consequences}
\label{sec:consequences}
\subsection{Centered regular tetrahedra}
Let $T(n)$ be the greatest common distance between the four vertices
of a regular tetrahedron contained in $C_n$, subject to the additional
condition that their centroid, the average of the four vertices,
is the origin.

This centering condition distinguishes our result from two related
problems. Hudelson, Klee, and Larman \cite{HudelsonKleeLarman1996}
maximize the volume of simplices in cubes without imposing regularity.
Firsching \cite{Firsching2017} discusses regular tetrahedra in cubes
without fixing the centroid, including constructions in dimensions
$4$ and $7$. Neither result supplies the centered optimum below.
We prove
\begin{equation}\label{eq:tetrahedron}
 T(n)=\sqrt2\,f(3,n).
\end{equation}
To see this, use the sign matrix
\[
 H=\begin{pmatrix}
 1&1&-1&-1\\
 1&-1&1&-1\\
 1&-1&-1&1
 \end{pmatrix}.
\]
It satisfies $HH^{\mathsf T}=4I_3$ and
$H^{\mathsf T}H=4I_4-J_4$. For a three-vector $a$,
$\norm{a^{\mathsf T}H}_\infty=\norm a_1$: the four columns of $H$
and their negatives are the eight sign vectors.
For a feasible matrix $A$ at level $s$, the four columns of $AH/2$
are centered, lie in $C_n$, and have pairwise distance $\sqrt{2s}$.
Conversely every centered regular tetrahedron of edge length $\ell$
has a vertex matrix $V$ with
$V^{\mathsf T}V=(\ell^2/2)(I_4-J_4/4)$. The matrix
$A=VH^{\mathsf T}/2$ then satisfies $AH/2=V$,
$A^{\mathsf T}A=(\ell^2/2)I_3$, and the required row bounds.
This proves \eqref{eq:tetrahedron}.

A shorter geometric argument is also available. Since $C_n$ is
convex and symmetric about the origin, a centered regular tetrahedron
$V$ lies in $C_n$ if and only if $\operatorname{conv}(V\cup -V)$
does, and $\operatorname{conv}(V\cup-V)$ is a centered three-cube
of side $\ell/\sqrt2$. Conversely the vertices of a centered
three-cube split into two such tetrahedra. Finally, centering costs
nothing in $f(3,n)$: if a cube $Q$ lies in $C_n$, so does $-Q$, and
the Minkowski average $\tfrac12(Q+(-Q))$ is a congruent centered cube.
In particular,
\[
 \F(5)=11-4\sqrt6,\qquad T(5)=4-\sqrt6.
\]
The equality concerns only tetrahedra with centroid at the
origin. It does not assert the same optimum without that
restriction.

\Needspace{12\baselineskip}
\subsection{General inner dimension}
For a fixed inner dimension $k\ge2$ and residue $0\le r<k$,
put
\[
 g_r(m)=F_k(km+r)-m\qquad(m\ge1).
\]
Subtracting $m$ removes the linear increase associated with
repeated stacking of identity blocks. The resulting \emph{residue
correction} stabilizes if there exists an integer $M_r$ such that
$g_r(m)=g_r(M_r)$ for every $m\ge M_r$.
Stacking and the trace bound give
\[
 g_r(m+1)\ge g_r(m),\qquad 0\le g_r(m)\le r/k
\]
whenever $m\ge1$. Thus $g_r(m)$ has a limit.  A bounded nondecreasing sequence need
not be eventually constant, so an additional structural condition is
required.

A \emph{unit coordinate row} is a signed standard basis vector
$\pm e_j$.  Such a row contributes one unit to the $j$th diagonal
entry of $A^{\mathsf T}A$.  Rows such as $(t,0,\ldots,0)$ with
$0<|t|<1$ are not removable unit coordinate rows, even though
their support has size one.  We call the remaining non-unit rows
the \emph{residual core}.

\Needspace{14\baselineskip}
\begin{proposition}[Finite cores and stabilization]\label{prop:finite-core}
Fix $k\ge2$. The following two properties are equivalent:
\begin{enumerate}
\item There is an integer $C$ such that, for every sufficiently large
ambient dimension $n$, some maximizing feasible matrix has at most $C$
rows that are not signed coordinate vectors.
\item For every residue $r$ modulo $k$, the sequence $g_r(m)$ is eventually
constant.
\end{enumerate}
\end{proposition}
\begin{proof}
Only signed unit coordinate rows are removed in this argument.
Suppose (1), and choose nonnegative integers $C,N$ such that the stated
core bound holds in every dimension $n\ge N$. Set
$M=\max\{1,C,N\}$. Fix $0\le r<k$ and $m\ge M$, and choose a maximizing
frame in dimension $km+r$ with at most $C$ residual rows. Its level is
$m+c$, where $c=g_r(m)$ and $0\le c\le r/k<1$.

Let $q_j$ be the number of unit coordinate rows in coordinate $j$, and let $t\le C$
be the number of residual rows. The energy in column $j$ gives
$q_j\le m+c<m+1$, hence $q_j\le m$. Counting all rows now yields
\[
 q_j=km+r-t-\sum_{\ell\ne j}q_\ell
 \ge m+r-t\ge m-C\ge m-M.
\]
Remove $m-M$ unit coordinate rows in each coordinate. These choices are disjoint,
and their total contribution to the Gram matrix is $(m-M)I_k$.
The remaining matrix has $kM+r$ rows, satisfies the same row bounds,
and has Gram matrix $(M+c)I_k$. Therefore
\[
 M+c\le F_k(kM+r),\qquad g_r(m)\le g_r(M).
\]
Stacking gives the reverse inequality, so $g_r(m)=g_r(M)$ for every
$m\ge M$. This proves (2).

Conversely, choose one maximizing frame at the start of the constant
tail in each residue class. Stack identity blocks to obtain maximizing
frames in every later dimension. The number of noncoordinate rows
never increases, so the finitely many chosen cores provide a uniform $C$.
\end{proof}

\begin{remark}
Removing one copy of $I_k$ reduces the number of rows by $k$
and the level by $1$.  The bound on residual rows ensures that
sufficiently many such blocks can be removed.
\end{remark}

The sharp formula in \cref{thm:sharp-values} is specific to
$k=3$. For general $k\ge2$, \cref{prop:finite-core} is an equivalence, not
an existence theorem for uniformly bounded residual cores.
In the case $k=3$, the constructions yield a concrete uniform
bound of four non-unit coordinate rows in some maximizer in every
dimension $n\ge3$.  In dimensions $3m$, the identity construction
has no such rows.  The four-row construction has at most four, the
seven-row construction has three, and the five- and ten-row cores
each have four. Stacking identity blocks preserves these counts.
This existence statement is also recorded as a separate
Lean theorem.  We do not assert that four is the smallest uniform
bound or classify all maximizing frames up to symmetry.

% END sections/consequences
\appendix
\section{Rational endpoint checks}\label{app:root-arithmetic}
We record the rational endpoint comparisons used in
\cref{lem:root-isolation}. The estimates use
$141421356/10^8<\sqrt2<141421357/10^8$.
The endpoint estimates, checked by squaring rational bounds, are
\par\smallskip
{\centering\small
\begin{tabular}{@{}c c c c@{}}
\toprule
test & bound on $C_j$ & bound on $S_j$ & conclusion\\
\midrule
$j=1$, $y=17/100$ & $C_1<236/1000$ & $S_1<947/1000$ & $C_1^2+S_1^2<1$\\
$j=2$, $y=347/2000$ & $C_2>7971/10000$ & $S_2>6047/10000$ & $C_2^2+S_2^2>1$\\
$j=2$, $y=7/40$ & $C_2<796/1000$ & $S_2<605/1000$ & $C_2^2+S_2^2<1$\\
\bottomrule
\end{tabular}\par}
\smallskip
For the first test, $10143/10000<x_1<10144/10000$.
For the second, $142481/100000<x_2<142482/100000$.
At the last endpoint $x_2=57/40$ exactly.


\clearpage
\section*{Formalization and AI disclosure}
The accompanying Lean~4 project formalizes the matrix formulation
and the main results. It provides version-pinned dependencies and
reproducible build and audit instructions. The available local build
and kernel-axiom audit records concern the formal statements and
proofs, and do not by themselves verify every sentence of the
mathematical exposition.

The following table lists, for the principal statements, the Lean
declarations that formalize them. A reader can compare each
declaration's type, printed by \texttt{KeyTheorems.lean}, with the
statement in the text.

\begin{center}\small
\begin{tabular}{@{}ll@{}}
\toprule
Statement & Lean declaration(s)\\
\midrule
\cref{thm:sharp-values} & \texttt{geometricValue\_eq\_sqrt\_formula},\\
 & \texttt{frameValue\_eq\_formula}\\
\cref{cor:eventual-recurrence} & \texttt{geometricValue\_squared\_recurrence\_iff}\\
\cref{prop:matrix-model} & \texttt{geometricSide\_iff\_feasibleValue},\\
 & \texttt{General.value\_trace\_equality\_iff}\\
\cref{prop:stacking} & \texttt{General.value\_superadditive}\\
\cref{lem:active-r1} & \texttt{active\_residue\_one}\\
\cref{lem:second-order} & \texttt{strict\_second\_order\_not\_local\_min}\\
\cref{prop:support-bound} & \texttt{full\_support\_bound\_at\_optimum}\\
\cref{lem:root-isolation} & \texttt{rootOneEquation\_unique},
 \texttt{rootTwoEquation\_unique}\\
\cref{lem:r2-all-active} & \texttt{active\_residue\_two}\\
\cref{thm:r1,thm:r2-sharp} & \texttt{canonical\_sharp\_values}\\
Equation \eqref{eq:tetrahedron} & \texttt{tetrahedronValue\_eq\_geometricValue}\\
\cref{prop:finite-core} & \texttt{General.finite\_core\_iff\_stabilization}\\
\bottomrule
\end{tabular}
\end{center}
All names are in the namespace \texttt{F3}. In
\texttt{active\_residue\_two} the hypothesis $\delta>1/5$ is explicit;
at the optimum it holds because $\Delta>1/5$.

OpenAI ChatGPT, including GPT-6 and Pro-mode assistance, was used
during 2026 to explore possible arguments, assist with literature
searches, draft and revise parts of the manuscript, and write and
debug parts of the Lean~4 code. Use of these tools does not by
itself establish mathematical correctness. A full author-side
mathematical review of the exposition and formal specifications has
not yet been completed, and the draft has not undergone peer review.

\medskip
\begin{thebibliography}{99}
\bibitem{Gardner1966}
M. Gardner,
Is it possible to visualize a four-dimensional figure?,
\emph{Scientific American} \textbf{215} (1966), no.~5, 138--143
(\emph{Mathematical Games}).

\bibitem{CFG1991}
H. T. Croft, K. J. Falconer, and R. K. Guy,
\emph{Unsolved Problems in Geometry},
Springer, New York, 1991, Problem~B4, p.~53.
\url{https://doi.org/10.1007/978-1-4612-0963-8}.

\bibitem{Shultz2013}
K. R. P. D. Shultz,
\emph{Largest $m$-cube in an $n$-cube: Partial solution},
KITP preprint NSF-ITP-13-142 (2013).
Notes written in 1996 and assembled in 2013 with a preface by G. Huber.
The work is cited in~\cite{HSW2018}.
\url{https://www.kitp.ucsb.edu/sites/default/files/preprints/2013/13-142.pdf}.

\bibitem{Finch2003}
S. R. Finch,
\emph{Mathematical Constants},
Cambridge University Press, Cambridge, 2003, Sec.~8.14, pp.~524--525.

\bibitem{Firsching2014}
M. Firsching,
Answer to ``Problem equivalent to `largest square in a cube',''
\emph{MathOverflow}, 19 October 2014 (updated 23 October 2014).
\url{https://mathoverflow.net/questions/15498/problem-equivalent-to-largest-square-in-a-cube}.

\bibitem{Schrijver1986}
A. Schrijver,
\emph{Theory of Linear and Integer Programming},
Wiley, Chichester, 1986.

\bibitem{HudelsonKleeLarman1996}
M. Hudelson, V. Klee, and D. G. Larman,
Largest $j$-simplices in $d$-cubes: Some relatives of the Hadamard
maximum determinant problem,
\emph{Linear Algebra Appl.} \textbf{241--243} (1996), 519--598.
\url{https://doi.org/10.1016/0024-3795(95)00541-2}.

\bibitem{Firsching2017}
M. Firsching,
Answer to ``Largest regular $k$-simplex inscribed in a $d$-cube,
$k<d$'',
\emph{MathOverflow}, 9 August 2017 (subsequently edited).
\url{https://mathoverflow.net/a/278340}.

\bibitem{HSW2018}
G. Huber, K. Pechenick Shultz, and J. E. Wetzel,
The $n$-cube is Rupert,
\emph{American Mathematical Monthly} \textbf{125} (2018), no.~6, 505--512.
\url{https://doi.org/10.1080/00029890.2018.1448197}.

\bibitem{SteiningerYurkevich2023}
J. Steininger and S. Yurkevich,
An algorithmic approach to Rupert's problem,
\emph{Mathematics of Computation} \textbf{92} (2023), 1905--1929.
\url{https://doi.org/10.1090/mcom/3831}.

\end{thebibliography}
\end{document}
